\chapter[\hspace{0.6em}时变处理的 G 方法]{时变处理的 G 方法}

上一章介绍了一个含时变处理与处理-混杂反馈的数据集。我们已经说明，把传统的混杂校正方法用于该数据集时，无法正确校正混杂：即使时变处理对结局的因果效应为零，传统校正方法给出的效应估计仍与零不同。

本章给出处理-混杂反馈存在时传统方法偏倚的解决方案：使用 g 方法——g 公式、IP 加权、g 估计，以及它们的双重稳健推广。仍用上一章的数据集，我们在此说明三种 g 方法都给出正确的（零）效应估计。对时不变处理，第 13 章已介绍 g 公式，第 12 章介绍边缘结构模型的 IP 加权，第 15 章介绍结构嵌套模型的 g 估计。本章在第 19 章所述的可识别性条件——序贯可交换性、正值性与一致性——下，逐一引入这三种用于比较静态处理策略的 g 方法。

\section{时变处理的 g 公式}

再考虑表 20.1 中序贯随机试验的数据；为方便起见，这里把它们重新列为表 21.1。假设我们只关心时不变处理 $A_1$ 的效应，即想对比平均反事实结局 $\E[Y^{a_1=1}]$ 与 $\E[Y^{a_1=0}]$。

\begin{table}[htbp]
\centering
\caption{}
\label{tab:21.1}
\small
\begin{tabular}{ccccc}
\toprule
$N$ & $A_0$ & $L_1$ & $A_1$ & 平均 $Y$\\
\midrule
2400 & 0 & 0 & 0 & 84\\
1600 & 0 & 0 & 1 & 84\\
2400 & 0 & 1 & 0 & 52\\
9600 & 0 & 1 & 1 & 52\\
4800 & 1 & 0 & 0 & 76\\
3200 & 1 & 0 & 1 & 76\\
1600 & 1 & 1 & 0 & 44\\
6400 & 1 & 1 & 1 & 44\\
\bottomrule
\end{tabular}
\end{table}

在本书第一、第二部分我们已经证明，在可识别性条件下，每个均值 $\E[Y^{a_1}]$ 都是条件均值 $\E[Y|A_1=a_1,L_1=l_1]$ 的加权平均，条件变量是时不变处理与混杂因子。具体而言，$\E[Y^{a_1}]$ 等于加权平均 $\sum_{l_1}\E[Y|A_1=a_1,L_1=l_1]f(l_1)$，其中 $f(l_1)=\Pr[L_1=l_1]$，因为如上一章所示，只需 $L_1$ 就能使受治疗者（$A_1=1$）与未受治疗者（$A_1=0$）条件可交换。这个加权平均就是 $\E[Y^{a_1}]$ 的 g 公式：把平均结局标准化到研究总体中混杂因子（这里只有 $L_1$）的分布。

但在表 21.1 的序贯随机试验中，处理 $\bar A=(A_0,A_1)$ 是时变的，且如上一章所见，存在处理-混杂反馈。这意味着不能指望传统校正方法无偏地估计时变处理 $\bar A$ 的因果效应。例如，即使在序贯可交换性成立的序贯随机试验中，传统方法也可能给不出“始终治疗”下的平均结局 $\E[Y^{a_0=1,a_1=1}]$ 与“从不治疗”下的平均结局 $\E[Y^{a_0=0,a_1=0}]$ 的有效估计。相反，g 公式可用于计算序贯随机试验中的反事实均值 $\E[Y^{a_0,a_1}]$。为此，需要推广上面时不变处理的 g 公式表达式。

在可识别性条件（第 19 章所述）下，$\E[Y^{a_0,a_1}]$ 的 g 公式仍是加权平均，但现在是条件均值 $\E[Y|A_0=a_0,A_1=a_1,L_1=l_1]$ 的加权平均，条件变量是为实现序贯可交换性所需的时变处理与混杂因子。权重是混杂因子 $L_1$ 在给定过去下的分布；在本例中，过去就是与干预相对应的过去处理取值。具体而言，g 公式
\[
\sum_{l_1}\E[Y|A_0=a_0,A_1=a_1,L_1=l_1]f(l_1|a_0)
\]
在（静态）序贯可交换性下等于 $\E[Y^{a_0,a_1}]$。也就是说，对时变处理，可识别性条件下反事实平均结局的 g 公式估计量，是把平均结局标准化到研究总体中混杂因子的分布，且表达式中的每个因子都以过去的处理与协变量历史为条件。在时不变情形下，处理与混杂因子都在单一时间点测量，并不需要对先前历史取条件。

\mnote{时变处理的 g 公式最早由 Robins (1986, 1987) 描述。}

只有当对任意满足 $f(l_1|a_0)\neq 0$ 的 $l_1$，总体中都存在 $(A_0=a_0,A_1=a_1,L_1=l_1)$ 的个体时，g 公式才是可计算的（即定义良好）。这等价于技术注记 19.2 给出的正值性定义，也是技术注记 3.1 中关于正值性的讨论在时变处理下的推广。

{\setlength{\emergencystretch}{1em}
让我们把 g 公式用于估计表 21.1 序贯随机试验中的因果效应 $\E[Y^{a_0=1,a_1=1}]-\E[Y^{a_0=0,a_1=0}]$。平均 $\E[Y^{a_0=0,a_1=0}]$ 的 g 公式估计为 $84\times0.25+52\times0.75=60$。平均 $\E[Y^{a_0=1,a_1=1}]$ 的 g 公式估计为 $76\times0.50+44\times0.50=60$。因此因果效应 $\E[Y^{a_0=1,a_1=1}]-\E[Y^{a_0=0,a_1=0}]$ 的估计为 0，与预期一致。传统方法失败之处，g 公式成功了。\par}

\mnote{在一个有 2 个时间点的研究中，“从不治疗”的 g 公式为 $\E[Y^{a_0=0,a_1=0}]=\E[Y|A_0=0,A_1=0,L_1=0]\times\Pr[L_1=0|A_0=0]+\E[Y|A_0=0,A_1=0,L_1=1]\times\Pr[L_1=1|A_0=0]$。}

\begin{figure}[htbp]
\centering
\begin{tikzpicture}[causal, font=\scriptsize, scale=0.95, every node/.style={transform shape,fill=white}]
  \node at (11.2,5.05) {$N$};
  \node at (12.5,5.05) {$\E[Y|A_1,L_1]$};
  \draw[black!55] (10.6,4.85) -- (13.2,4.85);
  \coordinate (root) at (0,0);
  \node[draw,circle,inner sep=0pt,minimum size=2.8mm] (a0) at (3.0,2.2) {};
  \node[draw,circle,inner sep=0pt,minimum size=2.8mm] (a1) at (3.0,-2.2) {};
  \draw[ed] (root) -- node[above,sloped,pos=0.55] {$A_0{=}0\ (.50)$} node[below,sloped,pos=0.55] {16{,}000} (a0);
  \draw[ed] (root) -- node[below,sloped,pos=0.55] {$A_0{=}1\ (.50)$} node[above,sloped,pos=0.55] {16{,}000} (a1);
  \node[draw,circle,inner sep=0pt,minimum size=2.8mm] (l00) at (6.2,3.30) {};
  \node[draw,circle,inner sep=0pt,minimum size=2.8mm] (l01) at (6.2,1.10) {};
  \node[draw,circle,inner sep=0pt,minimum size=2.8mm] (l10) at (6.2,-1.10) {};
  \node[draw,circle,inner sep=0pt,minimum size=2.8mm] (l11) at (6.2,-3.30) {};
  \draw[ed] (a0) -- node[above,sloped,pos=0.55] {$L_1{=}0\ (.25)$} node[below,sloped,pos=0.55] {4{,}000} (l00);
  \draw[ed] (a0) -- node[below,sloped,pos=0.55] {$L_1{=}1\ (.75)$} node[above,sloped,pos=0.55] {12{,}000} (l01);
  \draw[ed] (a1) -- node[above,sloped,pos=0.55] {$L_1{=}0\ (.50)$} node[below,sloped,pos=0.55] {8{,}000} (l10);
  \draw[ed] (a1) -- node[below,sloped,pos=0.55] {$L_1{=}1\ (.50)$} node[above,sloped,pos=0.55] {8{,}000} (l11);
  \node[draw,circle,inner sep=0pt,minimum size=2.2mm] (x000) at (9.4,3.85) {};
  \draw[ed] (l00) -- node[above,sloped,pos=0.55] {$A_1{=}0\ (.60)$} (x000);
  \node at (11.2,3.85) {2400};
  \node at (12.5,3.85) {84};
  \node[draw,circle,inner sep=0pt,minimum size=2.2mm] (x001) at (9.4,2.75) {};
  \draw[ed] (l00) -- node[above,sloped,pos=0.55] {$A_1{=}1\ (.40)$} (x001);
  \node at (11.2,2.75) {1600};
  \node at (12.5,2.75) {84};
  \node[draw,circle,inner sep=0pt,minimum size=2.2mm] (x010) at (9.4,1.65) {};
  \draw[ed] (l01) -- node[above,sloped,pos=0.55] {$A_1{=}0\ (.20)$} (x010);
  \node at (11.2,1.65) {2400};
  \node at (12.5,1.65) {52};
  \node[draw,circle,inner sep=0pt,minimum size=2.2mm] (x011) at (9.4,0.55) {};
  \draw[ed] (l01) -- node[above,sloped,pos=0.55] {$A_1{=}1\ (.80)$} (x011);
  \node at (11.2,0.55) {9600};
  \node at (12.5,0.55) {52};
  \node[draw,circle,inner sep=0pt,minimum size=2.2mm] (x100) at (9.4,-0.55) {};
  \draw[ed] (l10) -- node[above,sloped,pos=0.55] {$A_1{=}0\ (.60)$} (x100);
  \node at (11.2,-0.55) {4800};
  \node at (12.5,-0.55) {76};
  \node[draw,circle,inner sep=0pt,minimum size=2.2mm] (x101) at (9.4,-1.65) {};
  \draw[ed] (l10) -- node[above,sloped,pos=0.55] {$A_1{=}1\ (.40)$} (x101);
  \node at (11.2,-1.65) {3200};
  \node at (12.5,-1.65) {76};
  \node[draw,circle,inner sep=0pt,minimum size=2.2mm] (x110) at (9.4,-2.75) {};
  \draw[ed] (l11) -- node[above,sloped,pos=0.55] {$A_1{=}0\ (.20)$} (x110);
  \node at (11.2,-2.75) {1600};
  \node at (12.5,-2.75) {44};
  \node[draw,circle,inner sep=0pt,minimum size=2.2mm] (x111) at (9.4,-3.85) {};
  \draw[ed] (l11) -- node[above,sloped,pos=0.55] {$A_1{=}1\ (.80)$} (x111);
  \node at (11.2,-3.85) {6400};
  \node at (12.5,-3.85) {44};
\end{tikzpicture}
\caption{}
\label{fig:21.1}
\end{figure}

理解 g 公式的另一种方式是把它看作一次模拟。在 $Y$ 与 $\bar L$ 联合满足序贯可交换性时，g 公式模拟了若研究总体中每个人都遵循处理策略 $\bar a$ 时本可观测到的反事实结局 $Y^{\bar a}$ 与协变量历史 $\bar L^{\bar a}$。换言之，g 公式模拟（识别）了策略 $\bar a$ 下反事实 $(Y^{\bar a},\bar L^{\bar a})$ 的联合分布。要看清这一点，先看图 21.1 中因果解释的结构树图，它是表 21.1 数据的另一种表示。在上述可识别性条件下，g 公式可视为构造一棵新树的过程，新树中所有人都遵循策略 $\bar a$。例如，图 21.2 的因果解释结构树图显示了若所有人都遵循“始终治疗”策略（$a_0=1,a_1=1$）时本可观测到的反事实总体。

\begin{figure}[htbp]
\centering
\begin{tikzpicture}[causal, font=\scriptsize, scale=0.95, every node/.style={transform shape,fill=white}]
  \node at (11.2,5.05) {$N$};
  \node at (12.5,5.05) {$\E[Y|A_1,L_1]$};
  \draw[black!55] (10.6,4.85) -- (13.2,4.85);
  \coordinate (root) at (0,0);
  \node[draw,circle,inner sep=0pt,minimum size=2.8mm] (a1) at (3.0,0) {};
  \draw[ed] (root) -- node[above,sloped,pos=0.55] {$A_0{=}1\ (1)$} node[below,sloped,pos=0.55] {32{,}000} (a1);
  \node[draw,circle,inner sep=0pt,minimum size=2.8mm] (l0) at (6.2,1.8) {};
  \node[draw,circle,inner sep=0pt,minimum size=2.8mm] (l1) at (6.2,-1.8) {};
  \draw[ed] (a1) -- node[above,sloped,pos=0.55] {$L_1{=}0\ (.25)$} (l0);
  \draw[ed] (a1) -- node[below,sloped,pos=0.55] {$L_1{=}1\ (.75)$} (l1);
  \node[draw,circle,inner sep=0pt,minimum size=2.2mm] (x1) at (9.4,2.4) {};
  \node[draw,circle,inner sep=0pt,minimum size=2.2mm] (x2) at (9.4,1.2) {};
  \node[draw,circle,inner sep=0pt,minimum size=2.2mm] (x3) at (9.4,-1.2) {};
  \node[draw,circle,inner sep=0pt,minimum size=2.2mm] (x4) at (9.4,-2.4) {};
  \draw[ed] (l0) -- node[above,sloped,pos=0.55] {$A_1{=}0\ (0)$} (x1);
  \draw[ed] (l0) -- node[above,sloped,pos=0.55] {$A_1{=}1\ (1)$} (x2);
  \draw[ed] (l1) -- node[above,sloped,pos=0.55] {$A_1{=}0\ (0)$} (x3);
  \draw[ed] (l1) -- node[above,sloped,pos=0.55] {$A_1{=}1\ (1)$} (x4);
  \node at (11.2,2.4) {0};
  \node at (11.2,1.2) {16{,}000};
  \node at (12.5,1.2) {76};
  \node at (11.2,-1.2) {0};
  \node at (11.2,-2.4) {16{,}000};
  \node at (12.5,-2.4) {44};
\end{tikzpicture}
\caption{}
\label{fig:21.2}
\end{figure}

为模拟这个反事实总体，我们 (i) 在 $k=0$ 与 $k=1$ 时分别以概率 1 接受处理 $a_0=1$ 与 $a_1=1$；(ii) 赋予与原研究总体相同的概率 $\Pr[L_1=l_1|A_0=a_0]$ 与相同的均值 $\E[Y|A_0=a_0,A_1=a_1,L_1=l_1]$。

有两点很重要。第一，g 公式的取值取决于 $L$ 中纳入了什么（如果有的话）。举例来说，假设我们不收集 $L_1$ 的数据，因为我们错误地认为，本研究可用图 20.8 去掉 $L_1\to A_1$ 箭头后的因果图表示。于是我们认为 $L_1$ 不是混杂因子，因而不需要用于识别。那么在没有 $L_1$ 数据时，g 公式变成 $\E[Y|A_0=a_0,A_1=a_1]$，因为已无协变量历史可校正。然而，由于本研究实际由图 20.8 的因果图表示（其中处理分配 $A_1$ 受 $L_1$ 影响），不含 $L_1$ 的 g 公式不再具有因果解释。

第二，即便 g 公式具有因果解释，它的每个成分都可能没有因果解释。例如，考虑图 20.9 的因果图，其中只有静态序贯可交换性成立。包含 $L_1$ 的 g 公式正确识别了 $Y^{a}$ 的均值。

\mnote{在序贯可交换性下，$\Pr[L_1=l_1|A_0=a_0]=\Pr[L_1^{a_0}=l_1]$，且 $\E[Y|A_0=a_0,A_1=a_1,L_1=l_1]=\E[Y^{a_0,a_1}|L_1^{a_0}=l_1]$。于是 g 公式为 $\sum_{l_1}\E[Y^{a_0,a_1}|L_1^{a_0}=l_1]\Pr[L_1^{a_0}=l_1]$，它等于所需的 $\E[Y^{a_0,a_1}]$。}

值得注意的是，无论是否加上 $A_0$ 与 $A_1$ 指向 $Y$ 的箭头，g 公式都继续具有 $\E[Y^{\bar a}]$ 的因果解释，尽管它的两个成分——$\E[Y|A_0=a_0,A_1=a_1,L_1=l_1]$ 与 $\Pr[L_1=l_1|A_0=a_0]$——都完全没有因果解释。也就是说，$\Pr[L_1=l_1|A_0=a_0]\neq\Pr[L_1^{a_0}=l_1]$，且 $\E[Y|A_0=a_0,A_1=a_1,L_1=l_1]\neq\E[Y^{a_0,a_1}|L_1^{a_0}=l_1]$。在像图 20.1 与 20.2 所表示的序贯随机试验中，最后两个不等式将取等号。

\renewcommand{\FPname}{精要（Fine Point）}
\begin{finepoint}[处理与协变量历史]
在描述 g 方法时，我们常提到为实现序贯可交换性所需的处理与协变量历史。对 g 公式，我们说它的各成分以先前的处理与协变量历史为条件。例如，与 $k=2$ 时离散混杂因子 $L_2$ 的概率相对应的因子
\[
f(l_2|\bar a_1,\bar l_1)=\Pr[L_2=l_2|A_0=a_0,A_1=a_1,L_0=l_0,L_1=l_1]
\]
以时刻 0 与 1 的处理和混杂因子为条件；$k=3$ 时的因子以时刻 0、1、2 的处理和混杂因子为条件，依此类推。

然而，“历史”一词不必按时间定义，因为如精要 7.4 所述，混杂因子在理论上可以位于处理的时间未来。反过来，如随图 7.4 的讨论所述，校正处理时间过去的某些变量可能引入选择偏倚（称为 M 偏倚）。因此，在本书中，$k$ 时刻因果相关的“历史”应理解为：为使处理 $A_k$ 满足条件可交换性所需的处理与混杂因子集合。多数情况下，这样使用的历史对应于按时间顺序的历史。
\end{finepoint}
\renewcommand{\FPname}{精要}

现在把 g 公式推广到有多个时刻 $k$ 的高维情形。g 公式为
\[
\sum_{\bar l}\E[Y|\bar A=\bar a,\bar L=\bar l]\prod_{k=0}^{K}f(l_k|\bar a_{k-1},\bar l_{k-1}),
\]
其中求和遍历所有可能的 $\bar l$ 历史（$\bar l_{k-1}$ 是截至 $k-1$ 时刻的历史）。和 $\sum_{\bar l}$ 也可写作 $\sum_{l_K}\cdots\sum_{l_1}\sum_{l_0}$。在给定每个时刻 $k$ 的 $(\bar L_k,\bar A_k)$ 时 $Y^{\bar a}$ 满足序贯可交换性的条件下，该表达式等于处理策略 $\bar a$ 下的反事实均值 $\E[Y^{\bar a}]$。精要 21.1 给出了“历史”一词更细致的定义。

\mnote{技术注记 21.1 给出 g 公式更一般的表达式，可用于计算密度，而不只是均值。}

然而在实践中，若数据是高维的（在含多个混杂因子或多个时间点的观察性研究中通常如此），g 公式的各成分无法直接计算。$\E[Y|\bar A=\bar a,\bar L=\bar l]$ 与 $f(l_k|\bar a_{k-1},\bar l_{k-1})$ 都需要估计。例如，我们可以拟合线性回归模型估计随访结束时结局变量的条件均值 $\E[Y|\bar A=\bar a,\bar L=\bar l]$，并拟合 logistic 回归模型估计每个时刻 $k\neq0$ 离散混杂因子 $L_k$ 的分布（$L_0$ 的分布无需模型即可估计，如 13.3 节所述）。然后把这些模型的估计 $\hat\E[Y|\bar A=\bar a,\bar L=\bar l]$ 与 $\hat f(l_k|\bar a_{k-1},\bar l_{k-1})$ 代入 g 公式。自第 13 章起，我们把这一估计量称为\emph{插入式 g 公式}；当插入式 g 公式所用的估计基于参数模型时，我们称其为\emph{参数 g 公式}。

为简单起见，本章主要讨论确定性策略下的 g 公式。不过，在序贯可交换性下，g 公式也可用于计算随机处理策略 $f^{int}$ 下结局的反事实均值。随机（静态）策略的一个例子是“在每个时刻 $k$ 独立地以概率 0.3 治疗、以概率 0.7 不治疗”，此时 $f^{int}(1|\bar a_{k-1},\bar l_k)=0.3$。也就是说，$f^{int}(a_k|\bar a_{k-1},\bar l_k)$ 是处理策略（或干预）$f^{int}$ 下 $k$ 时刻接受处理 $a_k$ 的条件概率。

\renewcommand{\TPname}{技术注记（Technical Point）}
\begin{techpoint}[g 公式密度]
对确定性静态策略 $\bar a$，$(Y,\bar L)$ 的 g 公式密度在 $(y,\bar l)$ 处的取值为
\[
f(y|\bar a_K,\bar l_K)\prod_{k=0}^{K}f(l_k|\bar a_{k-1},\bar l_{k-1}).
\]
$Y$ 的静态 g 公式密度就是 $(Y,\bar L)$ 的 g 公式密度下 $Y$ 的边缘密度：
\[
\int\cdots\int f(y|\bar a_K,\bar l_K)\prod_{k=0}^{K}dF(l_k|\bar a_{k-1},\bar l_{k-1}),
\]
其中积分记号用于容纳 $L_k$ 的某些成分连续的情形。

对动态确定性策略 $g=(g_0,\dots,g_K)$（$g_k(\bar a_{k-1},\bar l_k)$ 取值于 $A_k$ 的支撑集），$(Y,\bar L)$ 与 $Y$ 的 g 公式密度只需把上式中的 $\bar a_k$ 换成 $\bar a_k^g$。这里 $a_k^g$ 由 $k=0,\dots,K$ 递归定义：$\bar a_k^g\equiv(g_0(a^g_{-1},l_0),\dots,g_k(\bar a^g_{k-1},l_k))$，其中 $\bar a^g_{-1}$ 定义为 0。当每个 $g_k(\bar a_{k-1},l_k)$ 都是常数函数时，静态策略是动态策略的特例。

更一般地，给定观测数据 $O=(\bar A,\bar X,Y)$ 与未观测数据 $U$，其中 $\bar X$ 是除处理 $\bar A$ 与结局 $Y$ 之外所有已测变量的集合，g 公式的输入是：(i) 一个确定性处理策略 $g$；(ii) 一个表示观测数据（及其未观测共同原因）的因果 DAG；(iii) $X$ 的一个子集 $\bar L$，我们希望对其校正；(iv) 与 DAG 拓扑一致的一种 $\bar L$、$\bar A$ 与 $Y$ 的全序，即每个变量都排在其祖先之后的顺序。向量 $L_k$ 由 $\bar L$ 中排于 $A_{k-1}$ 之后、$A_k$ 之前的所有变量组成。所选顺序通常（但不总是）是时间的，如精要 21.1 所讨论。当对所选顺序而言 $Y^g$ 的序贯可交换性与正值性都成立时，$Y$ 的 g 公式密度等于若研究总体中所有人都遵循策略 $g$ 时本可观测到的密度 $f_{Y^g}(y)$。否则，g 公式仍可计算，但缺少因果解释。当正值性与 $(Y^g,\bar L^g)$ 的可交换性成立时（例如，$U$ 或 $\bar X\setminus\bar L$ 中没有任何变量有箭头直接指向任何处理变量），$(Y,\bar L)$ 的 g 公式密度等于密度 $f_{Y^g,\bar L^g}(y,l)$。
\end{techpoint}
\renewcommand{\TPname}{技术注记}

\mnote{R 代码：\texttt{gfoRmula} R 包（Lin et al. 2019）可从 CRAN 获取；\texttt{GFORMULA} SAS 宏可从 GitHub 获取。见本书网站。}

那么，一般的 g 公式表达式为
\[
\sum_{\bar a,\bar l}\E[Y|\bar A=\bar a,\bar L=\bar l]\prod_{k=0}^{K}f(l_k|\bar a_{k-1},\bar l_{k-1})\prod_{k=0}^{K}f^{int}(a_k|\bar a_{k-1},\bar l_k).
\]
注意，若把 $f^{int}(a_k|\bar a_{k-1},\bar l_k)$ 换成观测到的处理条件概率 $f(a_k|\bar a_{k-1},\bar l_k)$，这就是观测到的 $Y$ 均值的公式。

g 公式的这一表达式足够一般，可容纳确定性与随机策略。在确定性处理策略下，对策略规定的 $a_k$ 取值，$f^{int}(a_k|\bar a_{k-1},\bar l_k)$ 恒为 1，对其他取值为 0。例如，在“从不治疗”策略即 $\bar a=(0,0,\dots,0)$ 下，所有 $k$ 都有 $f^{int}(0|\bar a_{k-1},\bar l_k)=1$。由于对策略规定的处理取值 $f^{int}(a_k|\bar a_{k-1},\bar l_k)$ 等于 1、对其他处理取值等于 0，上述公式中不必包含 $f^{int}$ 因子，也不必对 $\bar a$ 求和。我们公开可得的软件实现了这一一般形式的 g 公式，因而能容纳任何处理策略。

\section{时变处理的 IP 加权}

假设我们只关心表 21.1 中时不变处理 $A_1$ 的效应，于是想对比反事实平均结局 $\E[Y^{a_1=1}]$ 与 $\E[Y^{a_1=0}]$。如第 12 章所见，在可识别性条件下，每个反事实均值 $\E[Y^{a_1}]$ 都等于由个体特定非稳定权重 $W^{A_1}=1/f(A_1|L_1)$ 或稳定权重 $SW^{A_1}=f(A_1)/f(A_1|L_1)$ 所造伪总体中的均值 $\E_{ps}[Y|A_1=a_1]$。粗略地说，IP 权重的分母是一个人在其混杂因子取值条件下接受他实际所接受的处理取值的概率。可以由伪总体中 $A_1=a_1$ 的受试者 $Y$ 的平均值，从观测研究数据估计 $\E_{ps}[Y|A_1=a_1]$。

当处理与混杂因子是时变的，这些时不变处理的 IP 权重需要推广。对两个时间点的时变处理 $\bar A=(A_0,A_1)$ 与时变协变量 $\bar L=(L_0,L_1)$，非稳定 IP 权重为
\[
W^{\bar A}=\frac{1}{f(A_0|L_0)}\times\frac{1}{f(A_1|A_0,L_0,L_1)}=\prod_{k=0}^{1}\frac{1}{f(A_k|\bar A_{k-1},\bar L_k)},
\]
稳定 IP 权重为
\[
SW^{\bar A}=\frac{f(A_0)}{f(A_0|L_0)}\times\frac{f(A_1|A_0)}{f(A_1|A_0,L_0,L_1)}=\prod_{k=0}^{1}\frac{f(A_k|\bar A_{k-1})}{f(A_k|\bar A_{k-1},\bar L_k)},
\]
其中 $A_{-1}$ 按定义为 0。粗略地说，时变处理的 IP 权重分母是一个人在其处理与协变量历史条件下接受他实际所接受的处理史的概率。

假设我们想对比反事实均值 $\E[Y^{a_0=1,a_1=1}]$ 与 $\E[Y^{a_0=0,a_1=0}]$。在静态策略的可识别性假设下，每个反事实均值 $\E[Y^{a_0,a_1}]$ 都等于由非稳定权重 $W^{\bar A}$ 或稳定权重 $SW^{\bar A}$ 所造伪总体中的均值 $\E_{ps}[Y|A_0=a_0,A_1=a_1]$。也就是说，每个反事实均值的 IP 加权估计量，是伪总体中 $\bar A=(A_0,A_1)$ 的个体 $Y$ 的平均值。

让我们把 IP 加权用于表 21.1 的数据。图 21.3 的因果解释结构树图，是图 21.1 的树图加上若干列：非稳定 IP 权重 $W^{\bar A}$，以及每个处理与协变量历史在相应伪总体中的人数 $N_W$。伪总体的大小为 128{,}000，即原总体的 32{,}000 人乘以 4（静态策略的个数）。由于本研究中没有 $L_0$，IP 权重的分母简化为 $f(A_0)f(A_1|A_0,L_1)$。

\mnote{与时不变处理的技术注记 12.2 结果类似，$\E_{ps}[Y|A_0=a_0,A_1=a_1]$ 等于 $\E\big[W^{\bar A}YI(A_0=a_0,A_1=a_1)\big]/\E\big[W^{\bar A}I(A_0=a_0,A_1=a_1)\big]=\E\big[SW^{\bar A}YI(A_0=a_0,A_1=a_1)\big]/\E\big[SW^{\bar A}I(A_0=a_0,A_1=a_1)\big]$，无论序贯可交换性是否成立，对非稳定与稳定伪总体都如此。}

\begin{figure}[htbp]
\centering
\begin{tikzpicture}[causal, font=\scriptsize, scale=0.72, every node/.style={transform shape,fill=white}]
  \node at (11.2,5.05) {$N$};
  \node at (12.4,5.05) {$\E[Y|A_1,L_1]$};
  \node at (13.9,5.05) {$W^{\bar A}$};
  \node at (15.2,5.05) {$N_W$};
  \node at (16.6,5.05) {$SW^{\bar A}$};
  \node at (17.9,5.05) {$N_{SW}$};
  \draw[black!55] (10.6,4.85) -- (18.6,4.85);
  \coordinate (root) at (0,0);
  \node[draw,circle,inner sep=0pt,minimum size=2.8mm] (a0) at (3.0,2.2) {};
  \node[draw,circle,inner sep=0pt,minimum size=2.8mm] (a1) at (3.0,-2.2) {};
  \draw[ed] (root) -- node[above,sloped,pos=0.55] {$A_0{=}0\ (.50)$} node[below,sloped,pos=0.55] {16{,}000} (a0);
  \draw[ed] (root) -- node[below,sloped,pos=0.55] {$A_0{=}1\ (.50)$} node[above,sloped,pos=0.55] {16{,}000} (a1);
  \node[draw,circle,inner sep=0pt,minimum size=2.8mm] (l00) at (6.2,3.30) {};
  \node[draw,circle,inner sep=0pt,minimum size=2.8mm] (l01) at (6.2,1.10) {};
  \node[draw,circle,inner sep=0pt,minimum size=2.8mm] (l10) at (6.2,-1.10) {};
  \node[draw,circle,inner sep=0pt,minimum size=2.8mm] (l11) at (6.2,-3.30) {};
  \draw[ed] (a0) -- node[above,sloped,pos=0.55] {$L_1{=}0\ (.25)$} node[below,sloped,pos=0.55] {4{,}000} (l00);
  \draw[ed] (a0) -- node[below,sloped,pos=0.55] {$L_1{=}1\ (.75)$} node[above,sloped,pos=0.55] {12{,}000} (l01);
  \draw[ed] (a1) -- node[above,sloped,pos=0.55] {$L_1{=}0\ (.50)$} node[below,sloped,pos=0.55] {8{,}000} (l10);
  \draw[ed] (a1) -- node[below,sloped,pos=0.55] {$L_1{=}1\ (.50)$} node[above,sloped,pos=0.55] {8{,}000} (l11);
  \node[draw,circle,inner sep=0pt,minimum size=2.2mm] (x000) at (9.4,3.85) {};
  \draw[ed] (l00) -- node[above,sloped,pos=0.55] {$A_1{=}0\ (.60)$} (x000);
  \node at (11.2,3.85) {2400};
  \node at (12.4,3.85) {84};
  \node at (13.9,3.85) {$\tfrac{1}{(.5)(.6)}$};
  \node at (15.2,3.85) {8000};
  \node at (16.6,3.85) {$\tfrac{(.5)(.3)}{(.5)(.6)}$};
  \node at (17.9,3.85) {1200};
  \node[draw,circle,inner sep=0pt,minimum size=2.2mm] (x001) at (9.4,2.75) {};
  \draw[ed] (l00) -- node[above,sloped,pos=0.55] {$A_1{=}1\ (.40)$} (x001);
  \node at (11.2,2.75) {1600};
  \node at (12.4,2.75) {84};
  \node at (13.9,2.75) {$\tfrac{1}{(.5)(.4)}$};
  \node at (15.2,2.75) {8000};
  \node at (16.6,2.75) {$\tfrac{(.5)(.7)}{(.5)(.4)}$};
  \node at (17.9,2.75) {2800};
  \node[draw,circle,inner sep=0pt,minimum size=2.2mm] (x010) at (9.4,1.65) {};
  \draw[ed] (l01) -- node[above,sloped,pos=0.55] {$A_1{=}0\ (.20)$} (x010);
  \node at (11.2,1.65) {2400};
  \node at (12.4,1.65) {52};
  \node at (13.9,1.65) {$\tfrac{1}{(.5)(.2)}$};
  \node at (15.2,1.65) {24{,}000};
  \node at (16.6,1.65) {$\tfrac{(.5)(.3)}{(.5)(.2)}$};
  \node at (17.9,1.65) {3600};
  \node[draw,circle,inner sep=0pt,minimum size=2.2mm] (x011) at (9.4,0.55) {};
  \draw[ed] (l01) -- node[above,sloped,pos=0.55] {$A_1{=}1\ (.80)$} (x011);
  \node at (11.2,0.55) {9600};
  \node at (12.4,0.55) {52};
  \node at (13.9,0.55) {$\tfrac{1}{(.5)(.8)}$};
  \node at (15.2,0.55) {24{,}000};
  \node at (16.6,0.55) {$\tfrac{(.5)(.7)}{(.5)(.8)}$};
  \node at (17.9,0.55) {8400};
  \node[draw,circle,inner sep=0pt,minimum size=2.2mm] (x100) at (9.4,-0.55) {};
  \draw[ed] (l10) -- node[above,sloped,pos=0.55] {$A_1{=}0\ (.60)$} (x100);
  \node at (11.2,-0.55) {4800};
  \node at (12.4,-0.55) {76};
  \node at (13.9,-0.55) {$\tfrac{1}{(.5)(.6)}$};
  \node at (15.2,-0.55) {16{,}000};
  \node at (16.6,-0.55) {$\tfrac{(.5)(.4)}{(.5)(.6)}$};
  \node at (17.9,-0.55) {3200};
  \node[draw,circle,inner sep=0pt,minimum size=2.2mm] (x101) at (9.4,-1.65) {};
  \draw[ed] (l10) -- node[above,sloped,pos=0.55] {$A_1{=}1\ (.40)$} (x101);
  \node at (11.2,-1.65) {3200};
  \node at (12.4,-1.65) {76};
  \node at (13.9,-1.65) {$\tfrac{1}{(.5)(.4)}$};
  \node at (15.2,-1.65) {16{,}000};
  \node at (16.6,-1.65) {$\tfrac{(.5)(.6)}{(.5)(.4)}$};
  \node at (17.9,-1.65) {4800};
  \node[draw,circle,inner sep=0pt,minimum size=2.2mm] (x110) at (9.4,-2.75) {};
  \draw[ed] (l11) -- node[above,sloped,pos=0.55] {$A_1{=}0\ (.20)$} (x110);
  \node at (11.2,-2.75) {1600};
  \node at (12.4,-2.75) {44};
  \node at (13.9,-2.75) {$\tfrac{1}{(.5)(.2)}$};
  \node at (15.2,-2.75) {16{,}000};
  \node at (16.6,-2.75) {$\tfrac{(.5)(.4)}{(.5)(.2)}$};
  \node at (17.9,-2.75) {3200};
  \node[draw,circle,inner sep=0pt,minimum size=2.2mm] (x111) at (9.4,-3.85) {};
  \draw[ed] (l11) -- node[above,sloped,pos=0.55] {$A_1{=}1\ (.80)$} (x111);
  \node at (11.2,-3.85) {6400};
  \node at (12.4,-3.85) {44};
  \node at (13.9,-3.85) {$\tfrac{1}{(.5)(.8)}$};
  \node at (15.2,-3.85) {16{,}000};
  \node at (16.6,-3.85) {$\tfrac{(.5)(.6)}{(.5)(.8)}$};
  \node at (17.9,-3.85) {4800};
\end{tikzpicture}
\caption{}
\label{fig:21.3}
\end{figure}

反事实均值 $\E[Y^{a_0=0,a_1=0}]$ 的 IP 加权估计量，是伪总体中的均值 $\E_{ps}[Y|A_0=0,A_1=0]$；我们把它估计为伪总体中 $(A_0=0,A_1=0)$ 的 32{,}000 人的平均结局。由图 21.3 的树，估计为 $84\times\frac{8000}{32000}+52\times\frac{24000}{32000}=60$。同样，$\E[Y^{a_0=1,a_1=1}]$ 的 IP 加权估计也是 60。因此因果效应 $\E[Y^{a_0=1,a_1=1}]-\E[Y^{a_0=0,a_1=0}]$ 的估计为 0，与预期一致。IP 加权与 g 公式一样，在传统方法失败之处成功了。

\mnote{用图 21.3 中的稳定权重 $SW^{\bar A}$ 也得到同样的 0 估计（请自行验证）。不过，在非稳定伪总体中 $\Pr_{ps}[A_k=1|\bar A_{k-1},\bar L_k]=1/2$，而在稳定伪总体中 $\Pr_{ps}[A_k=1|\bar A_{k-1}]=1/2$。}

注意，我们基于 g 公式的非参数估计与基于 IP 加权的估计完全相等。这一相等与因果推断无关。也就是说，即便可识别性条件不成立——因而 g 公式与 IP 加权估计都没有因果解释——两种方法也会给出同一个数。

让我们把 IP 加权推广到有多个时刻 $k=0,1,\dots,K$ 的高维情形。非稳定 IP 权重的一般形式为
\[
W^{\bar A}=\prod_{k=0}^{K}\frac{1}{f(A_k|\bar A_{k-1},\bar L_k)},
\]
稳定 IP 权重的一般形式为
\[
SW^{\bar A}=\prod_{k=0}^{K}\frac{f(A_k|\bar A_{k-1})}{f(A_k|\bar A_{k-1},\bar L_k)}.
\]
当可识别性条件成立时，这些 IP 权重造出一个伪总体，其中 (i) $Y^{\bar a}$ 的均值与实际总体中相同，但 (ii) 如图 19.1 那样，每个时刻 $k$ 的随机化概率为常数 $1/2$（非稳定权重），或至多依赖于过去的处理史（稳定权重）。因此平均因果效应 $\E[Y^{\bar a}]-\E[Y^{\bar a'}]$ 等于 $\E_{ps}[Y|\bar A=\bar a]-\E_{ps}[Y|\bar A=\bar a']$，因为两个伪总体中都成立序贯无条件可交换性。

\mnote{正文的描述只考虑静态策略。关于动态策略的 IP 加权描述，见技术注记 21.2。}

在真正的序贯随机试验中，量 $f(A_k|\bar A_{k-1},\bar L_k)$ 由设计已知。因此我们可以用它们计算非稳定 IP 权重，而各反事实均值及其差值的估计都保证无偏。相反，在观察性研究中，$f(A_k|\bar A_{k-1},\bar L_k)$ 需要从数据估计。数据高维时，我们可以例如拟合 logistic 回归模型，估计每个时刻 $k$ 二值处理的条件概率 $\Pr[A_k=1|\bar A_{k-1},\bar L_k]$。这些模型的估计 $\hat f(A_k|\bar A_{k-1},\bar L_k)$\allowbreak 随后替代 $W^{\bar A}$ 中的 $f(A_k|\bar A_{k-1},\bar L_k)$。若 $\hat f(A_k|\bar A_{k-1},\bar L_k)$ 基于误设的 $\Pr[A_k=1|\bar A_{k-1},\bar L_k]$ logistic 模型，则 $Y^{\bar a}$ 与 $Y^{\bar a}-Y^{\bar a'}$ 的估计将有偏。对稳定权重 $SW^{\bar A}$，我们还须为分子取得 $\hat f(A_k|\bar A_{k-1})$ 的估计。即便这一估计基于误设模型，$\E[Y^{\bar a}]$ 与 $\E[Y^{\bar a}]-\E[Y^{\bar a'}]$ 的估计仍无偏，尽管稳定伪总体中的 $\hat f(a_k|\bar a_{k-1})$ 不再依概率收敛到观测数据密度 $f(a_k|\bar a_{k-1})$。

\mnote{实践中，常见做法是对 $\Pr[A_k=1|\bar A_{k-1},\bar L_k]$ 拟合单一模型，其中含时刻 $k$ 的函数（一个时变截距）作为协变量，可能还含与其他协变量的乘积项，而不是对每个时刻 $k$ 分别拟合模型。}

\begin{techpoint}[动态处理策略的 IP 加权]
考虑确定性动态策略 $g=(g_0,\dots,g_K)$，其中 $g_k\equiv g_k(\bar a_{k-1},\bar l_k)$。策略 $g$ 下结局 $Y$ 的 g 公式等于 $\E\big[YI(\bar A_K=\bar A_K^g)W^{\bar A}\big]$，其中 $\bar a_K^g$ 已在技术注记 21.1 中定义。此外，$\E\big[YI(\bar A_K=\bar A_K^g)W^{\bar A}\big]=\E_{ps}[Y|\bar A_K=\bar A_K^g]=\E[Y|\bar A_K=\bar A_K^g]$，其中 $\E_{ps}[Y|\bar A_K=\bar A_K^g]$ 是伪总体中遵循策略 $g$ 的成员 $Y$ 的均值。与静态策略不同，$\E\big[YI(\bar A_K=\bar A_K^g)SW^{\bar A}\big]/\allowbreak\E\big[I(\bar A_K=\bar A_K^g)SW^{\bar A}\big]$ 不等于 g 公式，因为 $SW^{\bar A}$ 的分子依赖于 $A$。因此稳定权重不能用于动态策略。对随机动态策略 $f^{int}$，g 公式等于
\[
\E\Big[Y\prod_{k=0}^{K}f^{int}(A_k|\bar A_{k-1},\bar L_k)W^{\bar A}\Big]=\E\Big[Y\prod_{k=0}^{K}\frac{f^{int}(A_k|\bar A_{k-1},\bar L_k)}{f(A_k|\bar A_{k-1},\bar L_k)}\Big].
\]
\end{techpoint}

假设我们得到 $\E[Y^{\bar a}]$ 的两个估计，一个用参数 g 公式，另一个用通过参数模型估计的 IP 权重，且两者之差超出抽样变异所能合理解释的范围（两估计之差的抽样变异可用自助法量化）。于是我们可以断定，g 公式所用的参数模型或 IP 加权所用的参数模型（或两者）误设。无论可识别性假设是否成立，这一结论总是正确的。其含义是：你应始终用两种方法估计 $\E[Y^{\bar a}]$，若估计相差很大（按某个预先设定的标准），就重新检查所有模型并在必要时修改它们。下一节我们描述双重稳健估计量如何帮助应对模型误设。

\mnote{即便两种参数方法的估计相似，也没有模型未误设的逻辑保证，因为它们可能以同一方向同时有偏。}

此外，如上一节所讨论，未知量 $\E[Y^{\bar a}]$ 的个数远超样本量并不罕见。因此我们需要指定一个模型，把许多策略的信息结合起来，以帮助估计某个给定的 $\E[Y^{\bar a}]$。例如，我们可以假设处理史 $\bar a$ 对平均结局的效应，随累积处理 $cum(\bar a)=\sum_{k=0}^{K}a_k$ 线性增长。这一假设被编码在\emph{边缘结构均值模型}
\[
\E[Y^{\bar a}]=\beta_0+\beta_1cum(\bar a)
\]
之中（对所有 $\bar a$），它是第 12 章所讨论时不变处理边缘结构均值模型的更一般版本。模型左边有 $2^K$ 个不同的未知量，每个对应 $2^K$ 个不同策略 $\bar a$ 之一；右边却只有 2 个未知参数 $\beta_0$ 与 $\beta_1$。参数 $\beta_1$ 度量时变处理 $\bar A$ 的平均因果效应。平均因果效应 $\E[Y^{\bar a}]-\E[Y^{\bar a=0}]$ 等于 $\beta_1\times cum(\bar a)$。

\mnote{这个边缘结构模型是不饱和的。记住，饱和模型左右两边的未知量个数相等。}

如第 12 章所讨论，为估计边缘结构模型的参数，我们可以在稳定或非稳定伪总体中用普通最小二乘拟合线性回归模型
\[
\E[Y|\bar A]=\theta_0+\theta_1cum(\bar A).
\]
这在数学上等价于在原研究总体中用加权最小二乘拟合同一线性模型，权重分别为 $SW^{\bar A}$ 或 $W^{\bar A}$（在实际数据分析中，这些权重用其估计替代）。在可识别性条件下，$\theta_1$ 的加权最小二乘估计依概率收敛于因果参数 $\beta_1$，而不是关联参数 $\theta_1$。

\mnote{在统计学课程中，通常证明：在 $\E[Y|\bar A]$ 的模型正确指定下，普通最小二乘与加权最小二乘估计都依概率收敛于关联参数 $\theta_1$。该证明假设权重只依赖于 $\bar A$。当权重如本例那样依赖于给定 $\bar A$ 后与 $Y$ 相关的协变量 $\bar L$ 时，加权回归不再依概率收敛于 $\theta_1$。}

又如第 12 章所讨论，$\hat\beta_1$ 的方差——因而对比 $\E[Y^{\bar a}]-\E[Y^{\bar a=0}]$ 的方差——可用非参数自助法估计，或计算其解析方差（这需要额外的统计分析与编程）。我们也可以用 $\hat\beta_1$ 的稳健方差估计量构造保守的 95\% 置信区间，多数统计软件包直接输出该估计。对非饱和边缘结构模型，用权重 $SW^{\bar A}$ 拟合时区间宽度通常比用 $W^{\bar A}$ 更窄，因此 $SW^{\bar A}$ 权重更可取。

当然，若边缘结构均值模型误设，即若平均反事实结局通过 $cum(\bar a)$ 之外的某个时变处理函数（例如仅最后 5 个月的累积处理 $\sum_{k=K-5}^{K}a_k$）依赖于处理策略，或非线性地（例如二次地）依赖于累积处理，则 $\E[Y^{\bar a}]$ 的估计将是错的。然而，若我们以权重 $SW^{\bar A}$ 或 $W^{\bar A}$ 拟合模型
\[
\E[Y|\bar A]=\theta_0+\theta_1cum(\bar A)+\theta_2cum_{-5}(\bar A)+\theta_3cum(\bar A)^2,
\]
则对联合假设 $\theta_2=\theta_3=0$ 作二自由度 Wald 检验，就是检验“我们的边缘结构模型正确指定”这一零假设。也就是说，边缘结构模型的 IP 加权不受技术注记 21.3 所述 g 零悖论的影响。实践中，人们可能选择纳入处理史 $\bar A$ 的不同概括作为协变量、并使用三次样条等灵活函数的边缘结构模型。

\mnote{当使用 $SW^{\bar A}$ 权重并用自助法估计方差时，这一检验对上述特定方向的误设通常有良好的检验功效。}

最后，如 12.5 节所讨论，我们可以用边缘结构模型考察 $L_0$ 中协变量的子集 $V$ 所造成的效应修饰。例如，对二值基线变量 $V$，我们会把边缘结构均值模型细化为
\[
\E[Y^{\bar a}|V]=\beta_0+\beta_1cum(\bar a)+\beta_2V+\beta_3cum(\bar a)V.
\]
该模型的参数可用加权最小二乘拟合普通线性回归模型
\[
\E[Y|\bar A,V]=\theta_0+\theta_1cum(\bar A)+\theta_2V+\theta_3Vcum(\bar A)
\]
来估计，权重为 IP 权重 $W^{\bar A}$，或更好的
\[
SW^{\bar A}(V)=\prod_{k=0}^{K}\frac{f(A_k|\bar A_{k-1},V)}{f(A_k|\bar A_{k-1},\bar L_k)}.
\]
在处理-混杂反馈存在时，$V$ 只能包含基线变量。若 $V$ 含有 $k>0$ 的 $L_k$ 成分，则即使处理在任何时刻对平均结局都没有效应，参数 $\theta_1$ 与 $\theta_3$ 也可能不为 0。

下面我们描述对任意策略 $g$ 的反事实均值 $\E[Y^g]$ 的双重稳健估计量。

\begin{techpoint}[g 零悖论]
使用参数 g 公式时，即便可识别性条件成立，模型误设也会导致 $\E[Y^{\bar a}]$ 的估计有偏。假设存在处理-混杂反馈，且“处理对 $Y$ 无效应”这一严格零假设为真，即
\[
Y^{\bar a}-Y^{\bar a'}=0\quad\text{以概率 1 对一切 }\bar a'\text{ 与 }\bar a\text{ 成立}.
\]
那么，对任何策略 $\bar a$，$\E[Y^{\bar a}]$ 的 g 公式取值都相同，尽管如第 20 章所讨论，$\E[Y|\bar A=\bar a,\bar L=\bar l]$ 与 $f(l_k|\bar a_{k-1},\bar l_{k-1})$ 都依赖于 $\bar a$。现在假设我们用基于不同（即变异独立）参数 $\theta$ 与 $\varphi$ 的标准非饱和参数模型 $\E[Y|\bar A=\bar a,\bar L=\bar l;\theta]$ 与 $f(l_k|\bar a_{k-1},\bar l_{k-1};\varphi)$ 来估计 g 公式的各成分。于是 Robins 与 Wasserman (1997) 证明：当 $L_k$ 含有任何离散成分时，这些模型不可能全都正确指定，因为 $\E[Y^{\bar a}]$ 的 g 公式估计值一般会依赖于 $\bar a$。其后果是，基于估计的 g 公式所做的推断，在理论上可能导致严格零假设被错误拒绝，即便在序贯随机试验中也是如此。这一现象称为时变处理下估计的 g 公式的零悖论。进一步讨论见 Cox 与 Wermuth (1999) 以及 McGrath 等 (2022)。所幸，g 零悖论在实践中并未妨碍参数 g 公式给出零效应估计，这大概是因为该悖论引起的偏倚与典型的随机变异相比很小。

相反，如第 12、14 章所述，边缘结构均值模型的 IP 加权与结构嵌套均值模型的 g 估计都不受零悖论影响。无论我们为处理选择何种函数形式，这些模型在严格零假设下都正确指定。例如，边缘结构均值模型 $\E[Y^{\bar a}]=\beta_0+\beta_1cum(\bar a)$ 在零假设下正确指定，因为此时 $\beta_1=0$，$\E[Y^{\bar a}]$ 不依赖于 $\bar a$ 的函数。又如 21.4 节所定义，任何结构嵌套均值模型 $\gamma_k(\bar a_{k-1},l_k,\beta)$ 在严格零假设下都正确指定：真参数值为 $\beta=0$，且无论过去处理与协变量历史的函数如何，都有 $\gamma_k(\bar a_{k-1},l_k,\beta)=0$。
\end{techpoint}

\section{时变处理的双重稳健估计量}

第二部分简要提到结合 IP 加权与 g 公式的双重稳健方法。我们知道，IP 加权需要处理 $A$ 在混杂因子 $L$ 条件下的正确模型，g 公式需要结局 $Y$ 在处理 $A$ 与混杂因子 $L$ 条件下的正确模型。双重稳健方法要求处理 $A$ 的模型或结局 $Y$ 的模型中有一个正确。只要两个模型中至少一个正确（而且无需知道是哪一个正确），双重稳健估计量就依概率收敛于因果效应。精要 13.2 描述了时不变处理 $A$ 对结局 $Y$ 的平均因果效应的双重稳健插入式估计量。本节先回顾一个略有不同的时不变处理双重稳健插入式估计量，再把它推广到时变处理。

\mnote{当观察性研究如常见情形那样有许多混杂因子、需要非饱和模型时，双重稳健估计量给了我们两次做对的机会。}

假设我们要在可交换性、正值性与一致性下，构造含许多混杂因子 $L$ 的设定中时不变二值处理 $A$ 对二值结局 $Y$ 的平均因果效应 $\E[Y^{a=1}]-\E[Y^{a=0}]$ 的双重稳健估计量。我们将按技术注记 13.2 与 13.3 已讨论的方式构造 $\E[Y^a]$ 的双重稳健估计量。$\E[Y^{a=1}]$ 与 $\E[Y^{a=0}]$ 的双重稳健估计量之差，就是平均因果效应的双重稳健估计量。我们对 $\E[Y^a]$ 的双重稳健程序将使用结局模型 $\E[Y|A=a,L=l]$ 的估计与 $\Pr[A=1|L]$ 的模型，然后把它们适当结合。程序分三步。

第一步，由处理模型计算预测值 $\hat f(a|L)\equiv\widehat{\Pr}[A=a|L]$。第二步，在 $A=a$ 的个体中，用最大似然拟合含 $\hat W^a=1/\hat f(a|L)$ 作为协变量的线性 logistic 模型 $b(a,L;\theta)$（例如 $b(a,L;\theta)=\mathrm{expit}(\theta_{a,0}+\theta_{a,1}L+\theta_{a,2}\hat W^a)$），计算预测值 $\hat\E[Y|A=a,L]=b(a,L;\hat\theta)$。第三步，把 $\E[Y^{a=1}]$ 与 $\E[Y^{a=0}]$ 估计为标准化均值 $\tilde\E\big[b(1,L;\hat\theta)\big]$ 与 $\tilde\E\big[b(0,L;\hat\theta)\big]$，其中 $\tilde\E$ 表示对全部个体（受治疗与未受治疗）取样本平均。差值 $\tilde\E\big[b(1,L;\hat\theta)\big]-\tilde\E\big[b(0,L;\hat\theta)\big]$ 就是因果效应 $\E[Y^{a=1}]-\E[Y^{a=0}]$ 的双重稳健估计量。也就是说，在可识别性条件下，只要处理模型正确或结局模型正确，该估计量就依概率收敛于平均因果效应。必须认识到，$L$ 取值相同的受治疗者与未受治疗者，其 $b(1,L;\hat\theta)=\mathrm{expit}\big(\hat\theta_{1,0}+\hat\theta_{1,1}L+\hat\theta_{1,2}/\hat f(a=1|L)\big)$ 也相同；他们的 $b(0,L;\hat\theta)=\mathrm{expit}\big(\hat\theta_{0,0}+\hat\theta_{0,1}L+\hat\theta_{0,2}/\hat f(a=0|L)\big)$ 同样相同。

\mnote{这个双重稳健估计量归功于 Bang 与 Robins (2005)，与一个更早的估计量（Robins, 2000）密切相关。该估计量是一个目标最小损失估计量（TMLE），在 van der Laan 与 Rubin (2006)、van der Laan 与 Gruber (2012) 后来引入的术语中也称目标最大似然估计量。}

现在把这个双重稳健估计量推广到我们关心两个处理策略 $\bar a$ 与 $\bar a'$ 下反事实均值 $\E[Y^{\bar a}]$ 与 $\E[Y^{\bar a'}]$ 之比较的时变处理设定。时变处理下估计 $\E[Y^{\bar a}]$ 的双重稳健程序，遵循与时不变处理下估计 $\E[Y^a]$ 相同的三步。不过，如我们将看到的，第二步稍复杂，因为它需要拟合序贯回归模型。为简化记号，我们演示如何在“始终治疗”策略下得到 $\E[Y^{\bar a}]$ 的双重稳健估计量，即 $\bar a=\mathbf{1}$，其中 $\mathbf{1}=\mathbf{1}_K$ 是 $K+1$ 个 1 组成的向量。

第一步需要拟合回归模型 $\pi_k(\bar L_k;\alpha)$，用以逼近 $\pi_k(\bar L_k)=\Pr[A_k=1|\bar A_{k-1}=\mathbf{1}_{k-1},\bar L_k]$，模型在所有人与所有时刻 $k$ 上合并拟合。个体只有在截至 $k-1$ 一直（连续）受治疗，即 $\bar A_{k-1}=\mathbf{1}_{k-1}$ 时，才对 $k$ 时刻模型的拟合有贡献。然后我们用该模型的预测值 $\pi_k(\bar L_k;\hat\alpha)$，为那些截至 $m$ 一直受治疗（$\bar A_m=\mathbf{1}_m$）的个体估计时变 IP 权重
\[
W^{A_m}=\prod_{k=0}^{m}\frac{1}{f(A_k|\bar A_{k-1},\bar L_k)},
\]
它等于 $W^{\mathbf{1}_m}=\prod_{k=0}^{m}\frac{1}{\pi_k(\bar L_k)}$。也就是说，对一个始终受治疗的个体（$\bar A_K=\mathbf{1}_K$），我们为每个时刻 $m$ 赋一个不同的权重 $W^{\mathbf{1}_m}$，而不是像上一节那样只在随访结束时赋单一权重 $W^{\mathbf{1}_K}$。例如，若我们拟合参数模型 $\pi_k(\bar L_k;\alpha)=\mathrm{expit}(\alpha_{0,k}+\alpha_2L_k)$ 以逼近 $\Pr[A_k=1|\bar A_{k-1}=\mathbf{1},\bar L_k]$，则在我们表 21.1 的两时间点例子（$K=1$）中，预测值为 $\hat\pi_1=\mathrm{expit}(\hat\alpha_{0,1}+\hat\alpha_2L_1)$ 与 $\hat\pi_0=\mathrm{expit}(\hat\alpha_{0,0}+\hat\alpha_2L_0)$（因为 $A_{-1}\equiv0$），它们分别是 $\widehat{\Pr}[A_1=1|A_0=1,\bar L_1]$ 与 $\widehat{\Pr}[A_0=1|L_0]$ 的估计。这里我们用了缩写 $\hat\pi_k$ 表示 $\pi_k(\bar L_k;\hat\alpha)$。然后我们为截至 $m$ 一直受治疗的个体计算时变 IP 权重估计 $\hat W^{\mathbf{1}_m}=\prod_{k=0}^{m}\frac{1}{\hat\pi_k}$。第一步到此结束。

第二步需要在每个时刻 $m$ 拟合一个单独的结局模型 $b_m(\bar L_m;\beta_m)$，从最后时刻 $K$ 开始，到 $m=0$ 结束。$m$ 时刻的回归模型只用截至 $m$ 一直受治疗的个体拟合，并纳入 $\hat W^{\mathbf{1}_m}=\hat W^{A_m}$ 作为协变量。$K$ 时刻模型的因变量是 $Y$。$m<K$ 时，$m$ 时刻模型的因变量是 $m+1$ 时刻模型拟合得到的预测结局，即 $\hat B_{m+1}=\hat b_{m+1}(\bar L_{m+1};\hat\beta_{m+1})$。当我们到达预测值 $\hat B_0=\hat b_0(\bar L_0;\hat\beta)$ 时，第二步完成。

\mnote{对二值 $Y$，可拟合 logistic 模型 $b_m(\bar L_m;\beta_m)\equiv\mathrm{expit}[\gamma_mX_m+\varsigma_m\hat W^{\mathbf{1},m}]$，$X_m$ 是 $\bar L_m$ 的向量函数，$\beta_m=(\gamma_m,\varsigma_m)$；$\hat B_K$ 不是整数，但保证落在 $[0,1]$ 内，可用作 logistic 模型的结局变量。对连续 $Y$，可拟合线性回归模型 $\gamma_mX_m+\varsigma_m\hat W^{\mathbf{1},m}$。}

第三步，我们计算 $\hat\E[Y^{\bar a=\mathbf{1}}]$ 的估计，即 $\hat B_0$ 对全部个体的样本平均。若 (i) 对所有 $m$，结局模型 $b_m(\bar L_m;\beta_m)$ 都正确指定，或 (ii) 对所有 $m$，处理模型 $\pi_k(\bar L_k;\alpha)$ 都正确指定，则 $\hat\E[Y^{\bar a=\mathbf{1}}]$ 对 $\E[Y^{\bar a=\mathbf{1}}]$ 是（渐近）无偏的。此时称 $\hat\E[Y^{\bar a=\mathbf{1}}]$ 是\emph{双重稳健}的。然而，$\hat\E[Y^{\bar a=\mathbf{1}}]$ 实际上是\emph{多重稳健}的，因为若对某个 $m\in\{0,1,\dots,K-1\}$，处理模型对时刻 0 至 $m$ 正确、结局模型对时刻 $m+1$ 至 $K$ 正确，它也对 $\E[Y^{\bar a=\mathbf{1}}]$ （渐近）无偏。我们把估计量的这一性质称为 K+2 稳健性。技术注记 21.4 与 21.5 解释这些稳健性质为何成立，并说明存在比 $\hat\E[Y^{\bar a=\mathbf{1}}]$ 稳健性更好的估计量。事实上，我们展示了一个具有 $2^{K+1}$ 稳健性的 $\E[Y^{\bar a=\mathbf{1}}]$ 估计量。

\begin{techpoint}[K+2 稳健的增广 IP 加权估计量]
我们考虑 $K=1$ 的情形，论证可推广到任意 $K$。g 公式 $\psi$ 的 ICE 插入式估计量为 $\hat\psi_{gfor}=P_n[\hat b_0(L_0)]$，其中 $P_n$ 表示样本平均，$\hat b_0(L_0)=\hat\E[\hat b_1(L_0,L_1)|A_0=1,L_0]$，$\hat b_1(L_0,L_1)=\hat\E[Y|L_0,A_0=1,L_1,A_1=1]$。$\psi$ 的 IP 加权估计量为 $\hat\psi_{IPW}=P_n[A_0A_1Y/(\hat\pi_0\hat\pi_1)]$，其中 $\hat\pi_0$ 与 $\hat\pi_1$ 分别是 $\pi_0=\Pr(A_0=1|L_0)$ 与 $\pi_1=\Pr(A_1=1|L_0,L_1,A_0=1)$ 的估计。Robins 等 (1994) 导出了 $\psi$ 的增广 IP 加权估计量 $\hat\psi_{TR}=P_n[\hat U_{TR}]$，其中
\[
\hat U_{TR}=\frac{A_0A_1Y}{\hat\pi_0\hat\pi_1}-\frac{A_0}{\hat\pi_0}\left\{\frac{A_1}{\hat\pi_1}-1\right\}\hat b_1(L_0,L_1)-\left\{\frac{A_0}{\hat\pi_0}-1\right\}\hat b_0(L_0).
\]
下面我们说明 $\hat\psi_{TR}$ 是三重（即 $K+2$）稳健的。第一，若 $\hat\pi_0$ 与 $\hat\pi_1$ 依概率收敛，则 $\hat\psi_{TR}$ 对 $\psi$ 依概率收敛（单重稳健），因为此时 $\hat U_{TR}$ 后两项的样本平均依概率收敛于 0，而第一项的样本平均恰为 $\hat\psi_{IPW}$。第二，$\hat\psi_{TR}$ 是双重稳健的，因为当 $\hat\E[Y|L_0,A_0=1,L_1,A_1=1]$ 与 $\hat\E[\hat b_1(L_0,L_1)|A_0=1,L_0]$ 分别依概率收敛于 $\E(Y|L_0,A_0=1,L_1,A_1=1)$ 与 $\E[b_1(L_0,L_1)|A_0=1,L_0]$ 时，$\hat\psi_{TR}$ 依概率收敛。这里 $\hat\E[\hat b_1(L_0,L_1)|A_0=1,L_0]$ 把与得到 $\hat b_0(L_0)$ 时相同的回归算法施于真实的 $b_1(L_0,L_1)$。为看清这一点，整理各项得
\[
\hat U_{TR}=\hat b_0(L_0)+\frac{A_0A_1}{\hat\pi_0\hat\pi_1}\big(Y-\hat b_1(L_0,L_1)\big)+\frac{A_0}{\hat\pi_0}\big(\hat b_1(L_0,L_1)-\hat b_0(L_0)\big).
\]
后两项的样本平均依概率收敛于 0，第一项的样本平均为 $\hat\psi_{gfor}$。第三，$\hat\psi_{TR}$ 是三重稳健的，因为若 $\hat b_1(L_0,L_1)$ 与 $\hat\pi_0$ 都依概率收敛，它也依概率收敛（Molina 等 2017）。这可由把 $\hat U_{TR}$ 改写为
\[
\frac{A_0\hat b_1(L_0,L_1)}{\hat\pi_0}+\frac{A_0A_1}{\hat\pi_0\hat\pi_1}\big(Y-\hat b_1(L_0,L_1)\big)-\left(\frac{A_0}{\hat\pi_0}-1\right)\hat b_0(L_0)
\]
看出。于是后两项的样本平均收敛于零，第一项的样本平均收敛于 $\E[b_0(L_0)]$。不过，当只有 $\hat\pi_1$ 与 $\hat\E[b_1(L_0,L_1)|A_0=a_0,L_0]$ 依概率收敛时，它并不依概率收敛。

修改估计量 $\hat\psi_{TR}$，可以构造四重稳健（即 $2K+1$ 稳健）估计量 $\hat\psi_{QR}$：当只有 $\hat\pi_1$ 与 $\hat\E[b_1(L_0,L_1)|A_0=a_0,L_0]$ 依概率收敛时，它依概率收敛（Tchetgen Tchetgen 2009）。令
\[
\tilde b_0(L_0)=\hat\E\left[\frac{A_1Y}{\hat\pi_1}-\left(\frac{A_1}{\hat\pi_1}-1\right)\hat b_1(L_0,L_1)\Big|A_0=1,L_0\right].
\]
则 $\hat\psi_{QR}=P_n[\hat U_{QR}]$，其中 $\hat U_{QR}$ 与 $\hat U_{TR}$ 相同，只是把 $\hat b_0(L_0)$ 换成 $\tilde b_0(L_0)$。$\tilde b_0(L_0)$ 优于 $\hat b_0(L_0)$ 之处在于：$\tilde b_0(L_0)$ 本身是双重稳健的，即当 $\hat\E[b_1(L_0,L_1)|A_0=1,L_0]$ 依概率收敛于 $b_0(L_0)$，且 $\hat\pi_1$ 或 $\hat b_1(L_0,L_1)=\hat\E[Y|L_0,A_0=1,L_1,A_1=1]$ 依概率收敛时，$\tilde b_0(L_0)$ 依概率收敛于 $b_0(L_0)=\E[b_1(L_0,L_1)|A_0=1,L_0]$；这意味着 $\hat\psi_{QR}$ 是四重稳健的。
\end{techpoint}

\begin{techpoint}[插入式 K+2 稳健估计量]
技术注记 21.4 中估计量 $\hat\psi_{TR}$ 的一个潜在缺点是：对二值 $Y$，$\hat\psi_{TR}$ 在给定样本中可能落在 $\psi$ 的支撑集 $[0,1]$ 之外。相反，$\hat\psi_{gfor}=P_n\big[\hat b_0(L_0)\big]$ 是 $\psi$ 的插入式估计量；若用（参数或非参数）logistic 回归模型估计 $\E[Y|L_0,A_0=a_0,L_1,A_1=a_1]$ 与 $b_0(L_0)=\E[b_1(L_0,L_1)|A_0=a_0,L_0]$，它总落在 $[0,1]$ 内。若
\[
\tilde U_{TR}=\hat b_0(L_0)+\frac{A_0A_1}{\hat\pi_0\hat\pi_1}\big(Y-\hat b_1(L_0,L_1)\big)+\frac{A_0}{\hat\pi_0}\big(\hat b_1(L_0,L_1)-\hat b_0(L_0)\big)
\]
中可保证 $P_n\big[\frac{A_0A_1}{\hat\pi_0\hat\pi_1}(Y-\hat b_1(L_0,L_1))\big]$ 与 $P_n\big[\frac{A_0}{\hat\pi_0}(\hat b_1(L_0,L_1)-\hat b_0(L_0))\big]$ 在每个样本中都为零，就得到一个也是三重稳健的插入式估计量。例如，在 $b_1(L_0,L_1)=\E[Y|L_0,A_0=1,L_1,A_1=1]$ 的线性 logistic 模型中，对 $A_0=A_1=1$ 的个体以 $Y$ 为因变量用最大似然加入单变量项 $\theta_1\frac{A_0A_1}{\hat\pi_0\hat\pi_1}$，即得 $P_n\big[\frac{A_0A_1}{\hat\pi_0\hat\pi_1}(Y-\hat b_1(L_0,L_1))\big]=0$。接着，在 $b_0(L_0)\equiv\E[b_1(L_0,L_1)|A_0=a_0,L_0]$ 的 logistic 模型中，对 $A_0=1$ 的个体以 $\hat b_1(L_0,L_1)$ 为因变量最大化 logistic 似然并加入项 $\theta_0\frac{A_0}{\hat\pi_0}$，即得 $P_n\big[\frac{A_0}{\hat\pi_0}(\hat b_1(L_0,L_1)-\hat b_0(L_0))\big]=0$。正文给出的估计量 $\hat\E[Y^{a=1}]$ 是 $\hat\psi_{TR,plug}$ 的一个实例。
\end{techpoint}

\mnote{Molina 等 (2017) 指出，该估计量实际上是 K+2 稳健的。Rotnitzky 等 (2017) 研究了在用非参数与机器学习方法估计处理回归函数与结局回归函数时，该估计量及其他多重稳健估计量的渐近偏倚。}

\begin{techpoint}[多重稳健估计量]
以 $f^g(a_m|\bar a_{m-1},\bar l_m)$ 表示策略 $g$ 下时刻 $m$ 的处理密度。对静态策略 $\bar a^*$，$f^g(a_m|\bar a_{m-1},\bar l_m)=I(a_m=a_m^*)$；对确定性动态策略 $g$，$f^g(a_m|\bar a_{m-1},\bar l_m)=I\big(a_m=g_m(\bar a_{m-1},\bar l_m)\big)$；对随机动态策略 $f^{int}$，$f^g(a_m|\bar a_{m-1},\bar l_m)=f^{int}(a_m|\bar a_{m-1},\bar l_m)$。再令 $C^g_k$ 在个体的观测处理史 $\bar A_k$ 与 $g$ 相容时等于 0，否则等于 1。下述算法计算 $\psi=\E[Y^g]$ 的多重稳健插入式估计量 $\hat\psi_{dr,plug}$，它基于 Rotnitzky 等 (2017) 提出的估计量，与 Robins (2000)、Bang 与 Robins (2005)、van der Laan 与 Gruber (2012)、Petersen 等 (2014) 的估计量密切相关。

\begin{enumerate}[label=\arabic*.,leftmargin=1.9em]
\item 对 $m=0,1,\dots,K$ 拟合模型 $f_m(a_m|\bar a_{m-1},\bar l_m;\alpha_m)$，用以逼近 $f^g(a_m|\bar a_{m-1},\bar l_m)$。得到参数向量 $\alpha_m$ 的 MLE $\hat\alpha_m$。对每个时刻 $m$，计算权重
\[
\hat W^{g,m}=\prod_{k=0}^{m}\frac{f^g(A_k|\bar A_{k-1},\bar L_k)}{f_k(A_k|\bar A_{k-1},\bar L_k;\hat\alpha_k)}.
\]
\item 令 $\hat T_{K+1}=Y$。
\item 递归地对 $m=K,K-1,\dots,0$：
\begin{enumerate}[label=(\alph*),leftmargin=1.7em]
\item 用迭代重加权最小二乘（IRLS），在 $C^g_m=0$ 的个体中拟合广义线性模型
\[
b_m(\bar A_m,\bar L_m;\gamma_m,\varsigma_m)=\phi\big[\gamma_md_m(\bar A_m,\bar L_m)+\varsigma_m\hat W^{g,m}\big],
\]
$\phi$ 为逆典型联系，用以逼近条件期望 $\E[\hat T_{m+1}|\bar A_m,\bar L_m,C^g_m=0]$；则 $(\hat\gamma_m,\hat\varsigma_m)$ 满足
\[
\tilde\E\left\{I(C^g_m=0)\begin{pmatrix}d_m(\bar A_m,\bar L_m)\\[2pt] \hat W^{g,m}\end{pmatrix}\Big(\hat T_{m+1}-b_m(\bar A_m,\bar L_m;\hat\gamma_m,\hat\varsigma_m)\Big)\right\}=0.
\]
\item 令 $\hat T_m=\sum_{a_m}b_m(a_m,\bar A_{m-1},\bar L_m;\hat\gamma_m,\hat\varsigma_m)f^g(a_m|\bar A_{m-1},\bar L_m)$。
\end{enumerate}
\item $\hat\psi_{dr,plug}=\tilde\E[\hat T_0]$。
\end{enumerate}

正如 Molina 等 (2017) 所指出的，$\hat\psi_{dr,plug}$ 是 K+2 稳健的：除双重稳健外，若对任意 $p\in\{1,\dots,K\}$，模型 $b_m(\bar A_m,\bar L_m;\gamma_m,\varsigma_m)$ 对 $m\in\{K,\dots,p\}$ 正确指定，且模型 $f_m(a_m|\bar a_{m-1},\bar l_m;\alpha_m)$ 对 $m\in\{p-1,\dots,0\}$ 正确指定，则它对 $\psi$ 也是（渐近）无偏的。

当 $\hat W^{g,m}$ 不作为协变量使用时，上述算法计算的是 g 公式 $\E[Y^g]$ 的迭代条件期望（ICE）估计量（精要 21.2），它是 g 公式的单重稳健估计量。
\end{techpoint}

为估计“从不治疗”策略下的反事实均值 $\E[Y^{\bar a=\mathbf{0}}]$，用 $\bar a=\mathbf{0}$ 重复上述步骤，其中 $\mathbf{0}$ 是 $K+1$ 个 0 组成的向量。两个策略下估计的均值之差，就是平均因果效应 $\E[Y^{\bar a=\mathbf{1}}]-\E[Y^{\bar a=\mathbf{0}}]$ 的多重稳健估计量。

这里描述的多重稳健估计量只能用于估计静态处理策略 $\bar a$ 下的反事实均值 $\E[Y^{\bar a}]$。技术注记 21.6 描述对处理策略 $g$（可以是静态或动态、确定性或随机）下反事实均值 $\E[Y^g]$ 的多重稳健估计量。该估计量有时称为目标最小损失估计量（TMLE）。

多重稳健估计量的实现历来受计算限制与缺乏易用软件的阻碍，尤其是基于风险率的生存分析。我们预期，在不久的将来会有相应软件，这些（用机器学习与样本拆分拟合的）多重稳健估计量在研究复杂处理策略对失效时间结局的效应时会更为常用。关于 g 公式的不同表示及其与上述估计量的联系，见精要 21.2。

\section{时变处理的 g 估计}

若我们只关心表 21.1 中时不变处理 $A_1$ 的效应，或许可以求助于第 14 章所述、用于时不变处理在协变量各层内条件因果效应的结构嵌套均值模型。那些模型只有一个方程，因为只有一个时间点 $k=0$。推广到时变处理，模型需要为数据中的每个时间点各列一个方程。对表 21.1 中两个时间点的时变处理 $\bar A=(A_0,A_1)$，我们指定一个（饱和的）加性结构嵌套均值模型，含两个方程
\[
\begin{aligned}
\text{时刻 }k=0:\quad & \E\big[Y^{a_0,a_1=0}-Y^{a_0=0,a_1=0}\big|A_0=a_0\big]=\beta_0a_0\\
\text{时刻 }k=1:\quad & \E\big[Y^{a_0,a_1}-Y^{a_0,a_1=0}\big|L_1^{a_0}=l_1,A_0=a_0,A_1^{a_0}=a_1\big]\\
&\qquad =a_1(\beta_{11}+\beta_{12}l_1+\beta_{13}a_0+\beta_{14}a_0l_1).
\end{aligned}
\]
由一致性，$k=1$ 的条件期望可写为 $\E\big[Y^{a_0,a_1}-Y^{a_0,a_1=0}\big|L_1=l_1,A_0=a_0,A_1=a_1\big]$。由于我们假设 $Y$ 满足序贯可交换性，我们可以且将会 (i) 把 $k=0$ 的条件期望换成 $\E\big[Y^{a_0,a_1=0}-Y^{a_0=0,a_1=0}\big]$，因为 $A_0=a_0$ 可以从条件事件中删去；(ii) 把 $k=1$ 的条件期望换成 $\E\big[Y^{a_0,a_1}-Y^{a_0,a_1=0}\big|L_1^{a_0}=l_1,A_0=a_0\big]$，因为 $A_1^{a_0}=a_1$ 可以从条件事件中删去。

因此，$k=1$ 时的方程在每个由 $(A_0,L_1)$ 定义的 4 个处理与协变量历史内，建模 $k=1$ 时处理的效应。模型的这一成分是饱和的，因为第二个方程中的 4 个参数 $\beta_1$ 在 4 个可能的过去处理与协变量历史水平内参数化了 $a_1$ 对 $Y$ 的效应。第一个方程建模 $k=1$ 时处理设为零时 $k=0$ 时处理的效应。模型的这一成分也是饱和的，因为它只有一个参数 $\beta_0$，用以估计唯一可能历史内的效应（没有先前的处理或协变量，所以每个人的历史都相同）。结构嵌套模型的两个方程，正是该模型被称为\emph{嵌套}的原因。第一个方程建模在时刻 0 受治疗、之后不再受治疗的效应；第二个方程建模在时刻 1 受治疗、之后不再受治疗的效应；若有更多时间点，依此类推。

\mnote{$a_1$ 的效应为：$A_0=0$、$L_1^{a_0=0}=0$ 者的 $\beta_{11}$；$A_0=0$、$L_1^{a_0=0}=1$ 者的 $\beta_{11}+\beta_{12}$；$A_0=1$、$L_1^{a_0=1}=0$ 者的 $\beta_{11}+\beta_{13}$；$A_0=1$、$L_1^{a_0=1}=1$ 者的 $\beta_{11}+\beta_{12}+\beta_{13}+\beta_{14}$。}

\mnote{由一致性，当 $A_0=a_0$ 时 $L_1^{a_0}=L_1$。}

\renewcommand{\FPname}{精要（Fine Point）}
\begin{finepoint}[g 公式的几种表示]
g 公式在数学上可以有几种表示。这些不同表示在非参数意义下等价，但在实践中导致不同的估计量。本书通篇强调的 g 公式表示，是标准化（流行病学行话）的推广形式。也就是说，对时不变处理，平均结局的 g 公式为 $\sum_l\E[Y|A=a,L=l]f(l)$；如本章所述，对时变处理为
\[
\sum_{\bar l}\E[Y|\bar A=\bar a,\bar L=\bar l]\prod_{k=0}^{K}f(l_k|\bar a_{k-1},\bar l_{k-1}).
\]
由于基于这一表示的插入式估计量需要估计混杂因子随时间的联合密度，我们称其为 g 公式的\emph{联合密度建模}估计量。

g 公式的另一种表示是迭代条件期望。对时不变处理，我们在 13.3 节隐含使用了这一 g 公式表示 $\E\big[\E[Y|A=a,L]\big]$。对时变处理，该表示是迭代条件期望（ICE），可以递归定义（Robins 1986）。基于 ICE 表示的插入式估计量需要拟合序贯预测算法（例如回归模型）。21.3 节与技术注记 21.4 描述了 ICE 估计量，我们在那里把它与 IP 权重的估计结合，构造双重（实际上是 K+2）稳健估计量。

g 公式的另一种表示是 IP 加权。事实上，如技术注记 2.3 对时不变处理所示，在正值性下标准化均值与 IP 加权均值相等。对时变处理同样如此（Robins and Rotnitzky, 1992; Robins, 1993; Young et al., 2014）。如本章所述，基于 g 公式的 IP 加权表示的估计量需要估计处理在给定过去处理与协变量历史下随时间的条件密度。我们把这类估计量称为 IP 加权估计量，而不是 g 公式估计量。
\end{finepoint}
\renewcommand{\FPname}{精要}

让我们用 g 估计来估计 $K=1$ 时结构嵌套模型的参数。我们采用与第 14 章相同的方法。先考虑每个个体 $i$ 的加性保秩结构嵌套模型
\[
\begin{aligned}
Y_i^{a_0,0}&=Y_i^{0,0}+\psi_{0}a_{0}\\
Y_i^{a_0,a_1}&=Y_i^{a_0,0}+\psi_{11}a_1+\psi_{12}a_1L_{1,i}^{a_0}+\psi_{13}a_1a_0+\psi_{14}a_1a_0L_{1,i}^{a_0},
\end{aligned}
\]
其中第二个方程限于 $A_0=a_0$ 的个体。也就是说，第二个方程实际是两个方程，一个用于 $A_0=1$ 的个体，一个用于 $A_0=0$ 的个体。这使我们能由一致性把 $L_{1,i}^{a_0}$ 换成 $L_{1,i}$；当如图 19.6 那样 $L_1$ 不满足序贯可交换性时，这一替换是从观测数据识别模型参数所必需的。为简化记号，我们把 $Y_i^{a_0=0,a_1=0}$ 记为 $Y_i^{0,0}$。

第一个方程是保秩模型，因为效应 $\psi_0$ 对每个个体完全相同。因此，若受试者 $i$ 的 $Y_i^{0,0}$ 大于受试者 $j$ 的 $Y_j^{0,0}$，则 $i$ 与 $j$ 在 $Y^{1,0}$ 上的排序不变——该模型在策略之间保持秩。此外，在第二个方程下，若受试者 $i$ 的 $Y_i^{1,0}$ 大于受试者 $j$ 的 $Y_j^{1,0}$，只有当两人的 $A_{0,i}$ 取值 $a_0$ 与 $L_{1,i}=L_{1,i}^{a_0}$ 取值 $l_1$ 都相同时，我们才能确定个体 $i$ 的 $Y_i^{1,1}$ 也大于个体 $j$ 的 $Y_j^{1,1}$。由于保秩以局部因素（即 $L_1^{a_0=1}$ 的取值）为条件，我们把第二个方程称为条件（或局部）保秩模型。

如第 14 章所讨论，由于未测的遗传与环境风险存在个体异质性，保秩在生物学上不可信。这就是我们主要关心结构嵌套均值模型的原因：它对是否存在由未测因素造成的个体间额外效应异质性完全不作假设。不过，在 $Y$ 满足序贯可交换性时，保秩模型的某一类 g 估计量（见下文）对均值模型的参数 $\beta$ 依概率收敛，即便保秩模型误设。

\mnote{证明见 Robins (1994)。注意，用 g 估计拟合非饱和结构嵌套均值模型并不需要正值性。}

g 估计的第一步是把模型与观测数据联系起来，如第 14 章对时不变处理所做。为此注意，由一致性，对恰好接受处理取值 $a_0$ 与 $a_1$ 的个体，反事实结局 $Y^{a_0,a_1}$ 等于观测结局 $Y$。形式上，对 $(A_0=a_0,A_1=a_1)$ 的个体，$Y^{a_0,a_1}=Y^{A_0,A_1}=Y$。同样，对 $(A_0=a_0,A_1=0)$ 的个体，$Y^{a_0,0}=Y^{A_0,0}$；对 $A_0=a_0$ 的个体，$L_1^{a_0}=L_1$。现在可以把结构嵌套模型用观测数据改写为
\[
\begin{aligned}
Y^{A_0,0}&=Y-(\psi_{11}A_1+\psi_{12}A_1L_1+\psi_{13}A_1A_0+\psi_{14}A_1A_0L_1)\\
Y^{0,0}&=Y^{A_0,0}-\psi_0A_0
\end{aligned}
\]
（为简化记号，我们省略个体下标 $i$）。

g 估计的第二步是用观测数据计算候选反事实 $H_1(\psi^\dagger)$ 与 $H_0(\psi^\dagger)$。为此，我们把结构嵌套模型中的参数真值 $\psi$ 换成某个值 $\psi^\dagger$：
\[
\begin{aligned}
H_1(\psi^\dagger)&=Y-\left(\psi^\dagger_{11}A_1+\psi^\dagger_{12}A_1L_1+\psi^\dagger_{13}A_1A_0+\psi^\dagger_{14}A_1A_0L_1\right)\\
H_0(\psi^\dagger)&=H_1(\psi^\dagger)-\psi^\dagger_0A_0.
\end{aligned}
\]
如第 14 章一样，目标是找到等于真值 $\psi$ 的参数值 $\psi^\dagger$。当 $\psi^\dagger=\psi$ 且 $\bar A_{k-1}=\bar a_{k-1}$ 时，候选反事实 $H_k(\psi^\dagger)$ 等于截至 $k-1$ 接受处理 $\bar a_{k-1}$、之后不接受处理的真实反事实 $Y^{\bar a_{k-1},0_k}$。现在我们可以用序贯可交换性在每个时间点进行 g 估计。精要 21.3 描述如何对我们饱和结构嵌套模型的参数 $\psi$ 做 g 估计。结果，结构嵌套模型的所有参数都为 0，这意味着任意静态或动态策略 $g$ 下的所有反事实均值 $\E[Y^g]$ 都等于 60。这一结果与 g 公式和 IP 加权所得一致。g 估计与 g 公式、IP 加权一样，在传统方法失败之处成功了。

\renewcommand{\FPname}{精要（Fine Point）}
\begin{finepoint}[用饱和结构嵌套模型做 g 估计]
$k=1$ 时的序贯可交换性意味着：在 $(A_0,L_1)$ 的四个联合层中任一层内，$A_1=1$ 者的 $Y^{A_0,0}$ 均值等于 $A_1=0$ 者的均值。因此，当 $\psi^\dagger=\psi$ 时，$H_1(\psi^\dagger)$ 的均值也必须相等。

先考虑层 $(A_0,L_1)=(0,0)$。由表 21.2 的数据行 1 与 2，我们看到 $H_1(\psi)$ 的均值在 $A_1=0$ 时为 84，在 $A_1=1$ 时为 $84-\psi_{11}$。因此 $\psi_{11}=0$。接着令对应层 $(A_0,L_1)=(0,1)$ 的数据行 3 与 4 中 $H_1(\psi)$ 的均值相等，得 $52=52-\psi_{11}-\psi_{12}$。既然 $\psi_{11}=0$，就断定 $\psi_{12}=0$。继续令数据行 5 与 6 中 $H_1(\psi)$ 的均值相等，得 $76=76-\psi_{11}-\psi_{13}$。既然 $\psi_{11}=\psi_{12}=0$，就断定 $\psi_{13}=0$。最后令数据行 7 与 8 中 $H_1(\psi)$ 的均值相等，得 $44=44-\psi_{11}-\psi_{12}-\psi_{13}-\psi_{14}$，故 $\psi_{14}=0$。

为估计 $\psi_0$，先把 $\psi_{11}$、$\psi_{12}$、$\psi_{13}$、$\psi_{14}$ 的取值代入表 21.2 中 $H_0(\psi)$ 均值的表达式。本例中所有参数都等于 0，故 $H_0(\psi)$ 的均值等于观测结局 $Y$ 的均值。然后我们用结构方程模型的第一个方程，从 $H_1(\psi)$ 的均值中减去 $\psi_0A_0$，算出表中每一数据行的 $H_0(\psi)$ 均值，如表 21.3 所示。$k=0$ 时的序贯可交换性 $Y^{0,0}\indep A_0$ 意味着：$A_0=1$ 的 16{,}000 名受试者与 $A_0=0$ 的 16{,}000 名受试者，其 $H_0(\psi)$ 均值相同。$A_0=0$ 者的 $H_0(\psi)$ 均值为 $84\times0.25+52\times0.75=60$；$A_0=1$ 者为 $(76-\psi_0)\times0.5+(44-\psi_0)\times0.5=60-\psi_0$。因此 $\psi_0=0$。g 估计到此完成。

\begin{center}
\begin{minipage}{0.92\textwidth}
\centering
\captionof{table}{}
\label{tab:21.2}
\small
\begin{tabular}{cccc}
\toprule
$A_0$ & $L_1$ & $A_1$ & $H_1(\psi)$ 均值\\
\midrule
0 & 0 & 0 & 84\\
0 & 0 & 1 & $84-\psi_{11}$\\
0 & 1 & 0 & 52\\
0 & 1 & 1 & $52-\psi_{11}-\psi_{12}$\\
1 & 0 & 0 & 76\\
1 & 0 & 1 & $76-\psi_{11}-\psi_{13}$\\
1 & 1 & 0 & 44\\
1 & 1 & 1 & $44-\psi_{11}-\psi_{12}-\psi_{13}-\psi_{14}$\\
\bottomrule
\end{tabular}

\end{minipage}
\end{center}

\begin{center}
\begin{minipage}{0.92\textwidth}
\centering
\captionof{table}{}
\label{tab:21.3}
\small
\begin{tabular}{cccc}
\toprule
$A_0$ & $L_1$ & $A_1$ & $H_0(\psi)$ 均值\\
\midrule
0 & 0 & 0 & 84\\
0 & 0 & 1 & $84-\psi_0$\\
0 & 1 & 0 & 52\\
0 & 1 & 1 & $52-\psi_0$\\
1 & 0 & 0 & $76-\psi_0$\\
1 & 0 & 1 & $76-\psi_0$\\
1 & 1 & 0 & $44-\psi_0$\\
1 & 1 & 1 & $44-\psi_0$\\
\bottomrule
\end{tabular}
\end{minipage}
\end{center}
\end{finepoint}
\renewcommand{\FPname}{精要}

\mnote{该末次脉冲函数满足 $\gamma_k(\bar a_{k-1},l_k,0)=0$，故在“处理无效应”的零假设下 $\beta=0$。}

然而在实践中，我们会遇到有多个时刻 $k$、每个时刻有多个协变量 $L_k$ 的观察性研究。一般地，结构嵌套均值模型的方程个数等于时间点数 $k=0,1,\dots,K$。我们在正文中讨论的结构嵌套均值模型的最一般形式如下（更一般的结构嵌套均值模型见技术注记 21.13）。对每个时刻 $k=0,1,\dots,K$，
\[
\begin{aligned}
&\E\left[Y^{\bar a_{k-1},a_k,0_{k+1}}-Y^{\bar a_{k-1},0_k}\Big|\bar L_k^{\bar a_{k-1}}=\bar l_k,\bar A_{k-1}=\bar a_{k-1},A_k=a_k\right]\\
&\qquad=a_k\gamma_k(\bar a_{k-1},l_k,\beta),
\end{aligned}
\]
其中 $(\bar a_{k-1},a_k,0_{k+1})$ 是一个静态策略，在时刻 0 到 $k-1$ 之间分配处理 $\bar a_{k-1}$，在时刻 $k$ 分配处理 $a_k$，从时刻 $k+1$ 到随访结束 $K$ 分配处理 0。策略 $(\bar a_{k-1},a_k,0_{k+1})$ 与 $(\bar a_{k-1},0_k)$ 的唯一区别是前者在 $k$ 时给处理 $a_k$，后者在 $k$ 时给处理 0。这里每个 $\gamma_k(\bar a_{k-1},l_k,\psi^\dagger)$ 是参数向量 $\psi^\dagger$ 的已知函数，当 $\psi^\dagger=0$ 时取 0；$\beta$ 是 $\psi^\dagger$ 的真值。

再次，在 $Y$ 的序贯可交换性下，我们可以从上述条件事件中删去 $A_k=a_k$。在 $K=1$ 的例子中，$\gamma_0(a_{-1},l_0,\beta)$ 就是 $\beta_0$（$l_0$ 与 $a_{-1}$ 都可取为恒等于 0），$\gamma_1(a_0,l_1,\beta)$ 就是 $\beta_{11}+\beta_{12}l_1+\beta_{13}a_0+\beta_{14}a_0l_1$。

于是，结构嵌套均值模型是关于 $k$ 时刻一个大小为 $a_k$ 的处理末次脉冲（last blip）对 $Y$ 均值之效应的模型，该效应是过去处理与协变量历史 $(\bar a_{k-1},l_k)$ 的函数 $\gamma_k(\bar a_{k-1},l_k,\beta)$。结构嵌套模型与边缘结构模型的关系见技术注记 21.7。

\begin{techpoint}[边缘结构模型与结构嵌套模型]
结构嵌套均值模型是半参数边缘结构均值模型，当且仅当对所有 $(\bar a_{k-1},\bar l_k,\beta)$，
\[
\gamma_k(\bar a_{k-1},\bar l_k,\beta)=\gamma_k(\bar a_{k-1},\beta)
\]
不依赖于 $\bar l_k$。具体而言，它是函数形式为
\[
\E[Y^{\bar a}]=\alpha_0+\sum_{k=0}^{K}a_k\gamma_k(\bar a_{k-1},\beta)
\]
的半参数边缘结构均值模型，其中 $\alpha_0=\E[Y^{\bar a=0}]$ 是未知常数。然而，这样的结构嵌套均值模型并不只是边缘结构均值模型，因为它还施加了额外的强假设：不存在由过去协变量历史造成的效应修饰。相反，边缘结构模型对是否存在时变协变量造成的效应修饰不作假设。

若我们指定结构嵌套均值模型 $\gamma_k(\bar a_{k-1},\beta)$，则可以用 g 估计或 IP 加权估计 $\beta$。然而，当结构嵌套均值模型（因而边缘结构均值模型）正确指定时，最有效的 g 估计量比最有效的 IP 加权估计量更有效，因为 g 估计利用“无过去协变量效应修饰”这一额外假设来提高效率。相反，假设边缘结构均值模型正确而结构嵌套均值模型不正确，即 $\gamma_k(\bar a_{k-1},\bar l_k,\beta)\neq\gamma_k(\bar a_{k-1},\beta)$，则 $\beta$ 与 $\E[Y^{\bar a}]$ 的 g 估计将有偏，而 IP 加权估计仍无偏。于是我们有一个经典的方差-偏倚权衡。给定边缘结构模型，若 $\gamma_k(\bar a_{k-1},\bar l_k,\beta)=\gamma_k(\bar a_{k-1},\beta)$，g 估计可提高效率；否则会引入偏倚。
\end{techpoint}

我们现在可以讨论含末次脉冲函数 $\gamma_k(\bar a_{k-1},l_k,\beta)$ 的一般结构嵌套均值模型的参数估计。为引出估计程序，我们将利用这样一个事实：参数真值为 $\psi$ 的正确指定局部保秩模型，也是参数真值 $\beta=\psi$ 的正确指定结构嵌套均值模型（反之不成立）。给定结构嵌套均值模型，我们可以定义
\[
H_k(\psi^\dagger)=Y-\sum_{j=k}^{K}A_j\gamma_j(\bar A_{j-1},\bar L_j,\psi^\dagger).
\]
参数向量真值为 $\psi$ 的正确指定局部保秩模型，等价于 $H_k(\psi)$ 恰好等于去除了时刻 $j$ 到 $K$ 的处理效应后的反事实 $Y^{\bar A_{k-1},0_k}$。特别地，$H_0(\psi)$ 是不接受任何处理时 $Y^0$ 的取值。然而，若局部保秩假设不正确（只要存在处理效应，几乎总是如此），但结构嵌套均值模型正确，我们仍有 $\E[H_k(\beta)|\bar A_k,\bar L_k]=\E[Y^{\bar A_{k-1},0_k}|\bar A_k,\bar L_k]$ 且 $\E[H_0(\beta)]=\E[Y^0]$。因此，若得到 $\hat\beta$ 的一致估计，就可以把 $\E[Y^0]$ 一致地估计为 $H_0(\hat\beta)$ 的样本平均。这正是 g 估计所提供的结果。

{\setlength{\emergencystretch}{1.5em}
有多个时间点或协变量时，我们需要拟合非饱和结构嵌套均值模型。例如，我们可能假设函数 $\gamma_k(\bar a_{k-1},l_k,\beta)$ 对所有 $k$ 都相同。最简单的模型是 $\gamma_k(\bar a_{k-1},l_k,\beta)=\beta_1$，它假设处理末次脉冲的效应对所有过去历史与所有时刻 $k$ 都相同。其他选项有 $\beta_1+\beta_2k$，假设效应随处理时刻 $k$ 线性变化；以及 $\beta_1+\beta_2k+\beta_3a_{k-1}+\beta_4l_k+\beta_5l_ka_{k-1}$，允许 $k$ 时刻处理的效应被最近的处理与协变量取值修饰。\par}

{\setlength{\emergencystretch}{3em}
为描述多时间点结构嵌套均值模型的 g 估计，假设非饱和模型为 $\gamma_k(\bar a_{k-1},l_k,\beta)=\beta_1$。\allowbreak 相应的保秩模型蕴含 $H_k(\psi^\dagger)=Y-\sum_{j=k}^{K}A_j\psi^\dagger$，\allowbreak 它对任意 $\psi^\dagger$ 都可由观测数据计算。我们随后选择远小于与远大于任何实质合理 $\psi$ 取值的 $\psi_{low}$ 与 $\psi_{up}$，并对网格 $\psi_{low},\psi_{low}+0.1,\psi_{low}+0.2,\dots,\psi_{up}$ 上的每个 $\psi^\dagger$，计算（对每个个体的每个时刻）$H_k(\psi^\dagger)$ 的值。\par}

然后，对每个 $\psi^\dagger$，我们拟合一个（跨时间合并的）logistic 回归模型
\[
\begin{aligned}
&\mathrm{logit}\Pr[A_k=1|H_k(\psi^\dagger),\bar L_k,\bar A_{k-1}]\\
&\qquad=\alpha_0+\alpha_1H_k(\psi^\dagger)+\alpha_2W_k
\end{aligned}
\]
以建模 $k=0,\dots,K$ 时接受处理的概率。这里 $W_k=w_k(\bar L_k,\bar A_{k-1})$ 是由个体的协变量与处理数据 $(\bar L_k,\bar A_{k-1})$ 算出的协变量向量，$\alpha_2$ 是未知参数的行向量，每人贡献 $K+1$ 个观测。$\beta$ 的 g 估计就是使 $\alpha_1$ 的估计最接近 0 的网格值 $\psi^\dagger$。我们可以通过定义 $\hat\beta$ 为使 $\alpha_1=0$ 的得分检验 p 值等于 1 的 $\psi^\dagger$ 值，来免去网格搜索。也就是说，$\hat\beta$ 是解
\[
\sum_{i=1,k=0}^{i=N,k=K}\{A_i-\mathrm{expit}(\hat\alpha_0+\hat\alpha_2W_{i,k})\}H_{i,k}(\psi^\dagger)=0
\]
的 $\psi^\dagger$ 值，其中 $\hat\alpha_0$ 与 $\hat\alpha_2$ 由把 $\alpha_1$ 项设为零后拟合上述 logistic 模型得到。可用标准方程求解器。若 (i) 结构嵌套均值模型正确，(ii) $Y$ 满足序贯可交换性，(iii) 模型 $\mathrm{logit}\Pr[A_k=1|\bar L_k,\bar A_{k-1}]=\alpha_0+\alpha_2W_k$ 正确，(iv) $H_k(\psi^\dagger)$ 线性进入上述 logistic 模型（即作为 $H_k(\psi^\dagger)$ 而非 $\{H_k(\psi^\dagger)\}^2$ 或任何其他非线性函数；见技术注记 14.2），则估计量 $\hat\beta$ 依概率收敛。

\mnote{检验 $\alpha_1=0$ 时 p 值大于 0.05 的 $\psi^\dagger$ 取值集合的上下限，就是 $\psi$ 的 95\% 置信区间的上下限。}

上述程序是第 14 章所述 g 估计程序到时变处理的推广。为简单起见，我们考虑了只有一个参数 $\beta_1$ 的结构嵌套模型，这意味着效应不随时刻 $k$ 变化，也不随处理与协变量历史变化。现在假设参数 $\beta$ 是向量。具体地，假设模型为 $\gamma_k(\bar a_{k-1},l_k,\beta)=\beta_0+\beta_1k+\beta_2a_{k-1}+\beta_3l_k+\beta_4l_ka_{k-1}$，于是 $\beta$ 是 5 维的，$l_m$ 是 1 维的。要估计 5 个参数，处理模型中需要 5 个额外的协变量。例如，我们可以拟合模型
\[
\begin{aligned}
&\mathrm{logit}\Pr[A_k=1|H_k(\psi^\dagger),\bar L_k,\bar A_{k-1}]\\
&\quad=\alpha_0+H_k(\psi^\dagger)(\alpha_1+\alpha_2k+\alpha_3A_{k-1}+\alpha_4L_k+\alpha_5L_kA_{k-1})\\
&\qquad+\alpha_6W_k.
\end{aligned}
\]
协变量的具体选择不影响 $\beta$ 点估计的一致性，但决定其置信区间的宽度。此前的 g 估计程序随后需要在 5 维网格上搜索，$\beta$ 的每个分量 $\beta_j$ 占一维。因此，若每个分量取 20 个网格点，5 维网格上就有 $20^5$ 个不同的 $\beta$ 值。然而，当 $\beta$ 的维数大于 2 时，用网格搜索找 g 估计 $\hat\beta$ 在计算上可能很困难。这种情况下，我们可以把 g 估计 $\hat\beta$ 定义为使 $\alpha_{1-5}=(\alpha_1,\dots,\alpha_5)^T=0$ 的得分检验 p 值等于 1 的 $\psi^\dagger$ 值，从而免去网格搜索。也就是说，$\hat\beta$ 是解如下 5 维估计方程的 $\psi^\dagger$ 值：
\[
\sum_{i=1,k=0}^{i=N,k=K}\{A_i-\mathrm{expit}(\hat\alpha_0+\hat\alpha_6^TW_{i,k})\}H_{i,k}(\psi^\dagger)(1,k,A_{i,k-1},L_{i,k},L_{i,k}A_{i,k-1})^T=0,
\]
其中 $\hat\alpha_0$ 与 $\hat\alpha_6$ 由把 $\alpha_{1-5}$ 设为零后拟合上述 logistic 模型得到。可用标准方程求解器。事实上，当结构嵌套均值模型对 $\beta$ 线性时（如本节讨论的所有例子），最后一个方程的解 $\hat\beta$ 有闭式。见技术注记 21.8，它还描述了估计量的多重稳健形式。

\mnote{$\beta_j$ 的 95\% 联合置信区间，是使 5 自由度得分检验在 5\% 水平上不拒绝的取值集合。一种计算量很大的做法是用单变量 95\% Wald 置信区间 $\beta_j\pm1.96$ 倍其标准误。}

\begin{techpoint}[线性结构嵌套均值模型的闭式估计量]
当 $\gamma_k(\bar a_{k-1},l_k,\beta)=\beta^TR_k$ 对 $\beta$ 线性（$R_k=r_k(\bar L_k,\bar A_{k-1})$ 是已知函数的向量，如我们讨论的所有例子），且给定模型 $\mathrm{logit}\Pr[A_k=1|\bar L_k,\bar A_{k-1}]=\alpha^TW_k$ 时，$\hat\beta$ 有显式闭式
\[
\hat\beta=\left\{\sum_{i=1,k=0}^{i=N,k=K}A_{i,k}X_{i,k}(\hat\alpha)Q_{i,k}S_{i,k}^{T}\right\}^{-1}\left\{\sum_{i=1,k=0}^{i=N,k=K}Y_iX_{i,k}(\hat\alpha)Q_{i,k}\right\},
\]
其中 $X_{i,k}(\alpha)=[A_{i,k}-\mathrm{expit}(\alpha^TW_{i,k})]$，$S_{i,k}=\sum_{i=1,j=k}^{i=N,j=K}R_{i,j}$，而维数为 $\beta$ 的函数 $Q_{i,k}=q_k(\bar L_{i,k},\bar A_{i,k-1})$ 的选择影响效率但不影响一致性。$Q_k$ 的最优选择见 Robins (1994)。

事实上，当 $\gamma_k(\bar a_{k-1},l_k,\beta)$ 对 $\beta$ 线性时，我们可以指定一个工作模型 $\varsigma^TD_k=\varsigma^Td_k(\bar L_k,\bar A_{k-1})$ 来逼近 $\E[H_k(\beta)|\bar L_k,\bar A_{k-1}]=\E[Y^{\bar A_{k-1},0_k}|\bar L_k,\bar A_{k-1}]$，并定义
\[
\begin{aligned}
\begin{pmatrix}\tilde\beta\\[2pt] \tilde\varsigma\end{pmatrix}
&=\left\{\sum_{i=1,k=0}^{i=N,k=K}
\begin{pmatrix}A_{i,k}X_{i,k}(\hat\alpha)Q_{i,k}\end{pmatrix}
\begin{pmatrix}S_{i,k}^{T} & D_{i,k}^{T}\end{pmatrix}\right\}^{-1}\\
&\qquad\times\left\{\sum_{i=1,k=0}^{i=N,k=K}Y_i
\begin{pmatrix}X_{i,k}(\hat\alpha)Q_{i,k}\\[2pt] D_{i,k}\end{pmatrix}\right\},
\end{aligned}
\]
从而得到 $\beta$ 的闭式 $2K+1$ 重稳健估计量 $\tilde\beta$。具体而言，若对每个 $k$，$\E[Y^{\bar A_{k-1},0_k}|\bar L_k,\bar A_{k-1}]$ 的模型 $\varsigma^TD_k$ 正确，或 $\mathrm{logit}\Pr[A_k=1|\bar L_k,\bar A_{k-1}]$ 的模型正确，则 $\tilde\beta$ 是 $\psi$ 的一致渐近正态估计量。
\end{techpoint}

给定结构嵌套均值模型参数的一致 g 估计 $\hat\beta$，最后一步是估计所关心策略 $g$ 下的反事实均值 $\E[Y^g]$。如前所述，$\E[Y^0]$ 可由样本平均 $\hat\E\big[H_0(\hat\beta)\big]$ 一致地估计。若不存在由过去协变量历史造成的效应修饰，即 $\gamma_k(\bar a_{k-1},l_k,\beta)=\gamma_k(\bar a_{k-1},\beta)$，则静态策略 $\bar a$ 下的 $\E[Y^{\bar a}]$ 估计为
\[
\hat\E[Y^{\bar a}]=\hat\E\big[Y^{0_K}\big]+\sum_{k=0}^{K}a_k\gamma_k(\bar a_{k-1},\tilde\beta).
\]
另一方面，若结构嵌套均值模型依赖于 $L_k$，或我们想估计动态策略 $g$ 下的 $\E[Y^g]$，则需要用技术注记 21.9 所述算法模拟 $L_k$。

\begin{techpoint}[g 估计结构嵌套均值模型后 {$\E[Y^g]$} 的估计]
假设可识别性假设成立，已得到结构嵌套均值模型 $\gamma_k(\bar a_{k-1},l_k,\beta)$ 的双重稳健 g 估计 $\tilde\beta$，并希望估计动态策略 $g$ 下的 $\E[Y^g]$。为此，可用如下蒙特卡罗算法的步骤：

\begin{enumerate}[label=\arabic*.,leftmargin=1.9em]
\item 把“始终不给治疗”时的平均响应 $\E\big[Y^{0_K}\big]$ 估计为 $H_0(\tilde\beta)$ 在 $N$ 名研究受试者上的样本平均。记该估计为 $\hat\E\big[Y^{0_K}\big]$。
\item 对数据拟合 $f(l_k|\bar a_{k-1},\bar l_{k-1})$ 的参数模型（跨人与时间合并），并令 $\hat f(l_k|\bar a_{k-1},\bar l_{k-1})$ 表示模型下 $f(l_k|\bar a_{k-1},\bar l_{k-1})$ 的估计。
\item 对 $v=1,\dots,V$：
\begin{enumerate}[label=(\alph*),leftmargin=1.7em]
\item 从 $\hat f(l_0)$ 抽取 $l_{v,0}$。
\item 对 $k=1,\dots,K$ 递归地从 $\hat f\big(l_k|\bar a_{v,k-1},\bar l_{v,k-1}\big)$ 抽取 $l_{v,k}$，其中 $\bar a_{v,k-1}=\bar g_{k-1}(\bar l_{v,k-1})$，即策略 $g$ 对应的处理史。
\item 令 $\hat\Delta_{g,v}=\sum_{j=0}^{j=K}a_{v,j}\gamma_j(\bar a_{v,j-1},\bar l_{v,j},\tilde\beta)$ 为 $Y^g-Y^{0_K}$ 的第 $v$ 个蒙特卡罗估计，其中 $a_{v,j}=g_j(\bar l_{v,j-1})$。
\end{enumerate}
\item 令 $\hat\E[Y^g]=\hat\E\big[Y^{0_K}\big]+\sum_{v=1}^{v=V}\hat\Delta_{g,v}/V$ 为 $\E[Y^g]$ 的估计。
\end{enumerate}

若 $f(l_k|\bar a_{k-1},\bar l_{k-1})$ 的模型、结构嵌套均值模型 $\gamma_k(\bar a_{k-1},l_k,\beta)$、处理模型 $\Pr[A_k=1|\bar L_k,\bar A_{k-1}]$ 或结局模型 $\E[Y^{\bar A_{k-1},0_k}|\bar L_k,\bar A_{k-1}]$ 中有一个正确指定，则 $\hat\E[Y^g]$ 是 $\E[Y^g]$ 的一致估计。置信区间可用非参数自助法得到。

注意，若估计 $\tilde\beta$ 对 $\beta=0$ 一致，则 $\gamma_k(\bar a_{k-1},l_k,\tilde\beta)$ 收敛于 0。因此，即便 $f(l_k|\bar a_{k-1},\bar l_{k-1})$ 的模型不正确，$\hat\Delta_{g,v}$ 也收敛于零，从而 $\hat\E[Y^g]$ 收敛于 $\hat\E\big[Y^{0_K}\big]$。也就是说，在可识别性条件成立、且我们对每个 $k$ 或者知道（如在序贯随机试验中）$\Pr[A_k=1|\bar L_k,\bar A_{k-1}]$、或者有 $\Pr[A_k=1|\bar L_k,\bar A_{k-1}]$ 或 $\E[Y^{\bar A_{k-1},0_k}|\bar L_k,\bar A_{k-1}]$ 的正确模型时，结构嵌套均值模型保持零假设。
\end{techpoint}

\section{删失是一种时变处理}

本章通篇使用一个没有删失的例子：表 21.1 中所有个体的结局都已知。然而在实践中，我们常会遇到某些个体失访、因而其结局值未知或（右）删失的情形。本书第二部分已讨论删失及其处理方法。第 8 章说明，即便在零假设下，删失也可能引入选择偏倚。第 12 章讨论了：我们通常关心的是若研究总体中无人被删失时的因果效应。

\mnote{关于删失，你可能想重读 12.6 节。}

然而在第二部分，我们只考虑了大大简化的删失，没有指明个体在随访期间何时被删失。也就是说，我们把删失 $C$ 当作时不变变量。更现实的看法是把删失视为时变变量 $C_1,C_2,\dots,C_{K+1}$，其中 $C_m$ 是指示变量，若个体在时刻 $m$ 仍未删失则取 0，否则取 1。删失是单调型缺失数据，即若个体 $C_m=0$，则此前所有删失指示也都为零（$C_1=0,C_2=0,\dots,C_{m-1}=0$）。此外，按定义，研究中所有个体都有 $C_0=0$；否则他们不会被纳入研究。

若个体在时刻 $m$ 被删失，即 $C_m=1$，则时刻 $m$ 之后测量的处理、混杂因子与结局都不可观测。因此，分析必然限于未删失的个体-时间，即 $C_m=0$ 者。例如，21.1 节反事实平均结局 $\E[Y^{\bar a}]$ 的 g 公式需要改写为
\[
\sum_{\bar l}\E[Y|\bar C=\bar 0,\bar A=\bar a,\bar L=\bar l]\prod_{k=0}^{K}f(l_k|c_k=0,\bar a_{k-1},\bar l_{k-1}),
\]
其中所有项都以仍未删失为条件。

\mnote{以未删失（$C=0$）为条件，在零假设下会引入选择偏倚：只要 $C$ 是处理 $A$ 与结局 $Y$ 之间某条路径上的对撞节点，或是这样一个对撞节点的后代。}

假设把每个时刻 $m$ 的处理 $A_m$ 换成 $(A_m,C_{m+1})$ 后可识别性条件成立。那么容易证明，上述表达式对应于联合处理 $(\bar a,\bar c=\bar 0)$ 下反事实平均结局 $\E[Y^{\bar a,\bar c=\bar 0}]$ 的 g 公式，即若所有人都接受处理策略 $\bar a$ 且无人失访时本可观测到的平均结局。

\mnote{上标 $\bar c=\bar 0$ 的使用，明确了很多人心中所指的处理 $\bar a$ 的因果对比，即便他们选择不使用上标 $c=0$。}

当联合处理 $(\bar A,\bar C)$ 的可识别性条件成立时，反事实均值 $\E[Y^{\bar a,\bar c=\bar 0}]$ 也可以通过结构均值模型的 IP 加权来估计。为估计该均值，我们可以例如把结局回归模型
\[
\E[Y|\bar A,\bar C=\bar 0]=\theta_0+\theta_1cum(\bar A)
\]
拟合到由非稳定 IP 权重 $W^{\bar A}\times W^{\bar C}$ 造出的伪总体，其中
\[
W^{\bar C}=\prod_{k=1}^{K+1}\frac{1}{\Pr(C_k=0|C_{k-1}=0,\bar A_{k-1},\bar L_{k-1})}.
\]
我们通过拟合 $\Pr(C_k=0|C_{k-1}=0,\bar A_{k-1},\bar L_{k-1})$ 的 logistic 回归模型来估计权重的分母。技术注记 21.10 说明如何推广到失效时间结局的生存分析。

在非稳定 IP 权重造出的伪总体中，被删失的个体由处理与协变量历史相同的未删失个体的副本替代。因此伪总体的大小与原研究总体在删失前（即任何失访发生前）相同。非稳定 IP 权重在伪总体中消除了删失。

或者我们可以用稳定 IP 权重 $SW^{\bar A}\times SW^{\bar C}$ 造出的伪总体，其中
\[
SW^{\bar C}=\prod_{k=1}^{K+1}\frac{\Pr(C_k=0|C_{k-1}=0,\bar A_{k-1})}{\Pr(C_k=0|C_{k-1}=0,\bar A_{k-1},\bar L_{k-1})}.
\]
我们通过分别拟合 $\Pr(C_k=0|C_{k-1}=0,\bar A_{k-1},\bar L_{k-1})$ 与 $\Pr(C_k=0|C_{k-1}=0,\bar A_{k-1})$ 两个模型，估计 IP 权重的分母与分子。

\mnote{记住：当 $\Pr(C_k=0|\bar A_{k-1},C_{k-1}=0,\bar L_k)$ 的模型正确指定时，估计的 IP 权重 $SW^{C}$ 均值为 1。}

稳定 IP 权重造出的伪总体，大小与删失后的原研究总体相同，即在每个时刻 $k$，伪总体中未删失个体的比例与研究总体相同。稳定权重并不在伪总体中消除删失，而是使删失在每个时刻 $k$ 相对于已测协变量历史 $\bar L_k$ 随机发生。也就是说，存在选择，但没有选择偏倚。无论使用哪类 IP 权重，伪总体中都没有从 $L_k$ 与 $A_k$ 指向未来 $C_m$（$m>k$）的箭头。重要的是，在联合处理 $(\bar A,\bar C)$ 的可交换性条件下，即便 $\bar L$ 的某些成分受先前处理影响，IP 加权也能无偏地估计 $(\bar A,\bar C)$ 的联合效应。

最后，使用结构嵌套模型的 g 估计时，我们需要先用 IP 加权校正删失所致的选择偏倚。实践中，这意味着我们先估计用于删失的非稳定 IP 权重 $W^{\bar C}$，造出一个无人被删失的伪总体，再对该伪总体做 g 估计。

\begin{techpoint}[时变处理的生存分析]
第 17 章描述了估计时点干预对失效时间结局之效应的 g 方法。本章描述估计持续策略对非失效时间结局之效应的 g 方法。实践中，我们常把第 17 章与本章的方法结合起来，用 g 方法估计持续策略对失效时间结局的效应。下面我们概述两种做法（分别基于 g 公式与 IP 加权），用以在序贯可交换性、正值性与一致性成立时估计处理策略 $\bar a$ 下的反事实风险 $\Pr\big[D^{\bar a,\bar c=\bar 0}_{k+1}=1\big]$。图 21.4 的因果图描绘了这样的设定：两个时间点，失效时间结局如第 17 章那样用时变指示变量表示。每个指示变量 $D_k$ 都应有箭头指向图上所有未来变量，但为减少杂乱，我们省略了这些箭头。为简单起见，我们也省略了删失的时变指示变量。

风险 $\Pr\big[D^{\bar a,\bar c=\bar 0}_{k+1}=1\big]$ 由 $1$ 减去 $\Pr\big[D^{\bar a,\bar c=\bar 0}_{k+1}=0\big]$ 的 g 公式来识别：
\[
\begin{aligned}
&\sum_{\bar l_k}\Pr\big[D_{k+1}=0|\bar A_k=\bar a_k,\bar L_k=\bar l_k,D_k=C_{k+1}=0\big]\\
&\quad\times\prod_{m=0}^{k}\Big\{f(l_m|\bar a_{m-1},\bar l_{m-1},D_m=C_m=0)\\
&\qquad\times\Pr\big[D_m=0|\bar A_{m-1}=\bar a_{m-1},\bar L_{m-1}=\bar l_{m-1},D_{m-1}=C_m=0\big]\Big\}.
\end{aligned}
\]
随后可以通过拟合离散时间风险率 $\Pr\big[D_{k+1}=1|\bar A_k=\bar a_k,\bar L_k=\bar l_k,D_k=C_{k+1}=0\big]$ 的模型，以及混杂因子 $L$ 随时间的条件密度 $f(l_k|\bar a_{k-1},\bar l_{k-1},D_k=C_k=0)$ 的模型，得到插入式 g 公式估计。如第 17 章所述，可用合并 logistic 模型逼近风险率。细节与应用见 Young 等 (2011)。Wen 等 (2021) 描述了 ICE g 公式估计量。

另一种做法是拟合风险率 $\Pr\big[D_{k+1}=1|\bar A_k=\bar a_k,\bar L_k=\bar l_k,D_k=C_{k+1}=0\big]$ 的合并 logistic 模型，其中每个个体在时刻 $k$ 获得时变非稳定 IP 权重 $W_k^{\bar A}\times W_k^{\bar C}$，而
\[
\begin{aligned}
W_k^{\bar A}&=\prod_{m=0}^{k}\frac{1}{f(A_m|\bar A_{m-1},D_m=C_m=0,\bar L_m)},\\
W_k^{\bar C}&=\prod_{m=1}^{k}\frac{1}{\Pr(C_m=0|\bar A_{m-1},D_{m-1}=C_{m-1}=0,\bar L_{m-1})},
\end{aligned}
\]
或在每个时刻 $k$ 使用相应的稳定 IP 权重。该模型的参数估计的是边缘结构合并 logistic 模型 $\Pr\big[D^{\bar a,\bar c=\bar 0}_{k+1}=1|D^{\bar a,\bar c=\bar 0}_k=0\big]$ 的参数（Robins 1998a）。细节与应用见 Hernán 等 (2001)。Wen 等 (2022) 综述了时变处理生存分析的多重稳健估计量。
\end{techpoint}

\begin{figure}[htbp]
\centering
\begin{tikzpicture}[causal, font=\small, scale=0.95, every node/.style={transform shape,fill=white}]
  \node[vtx] (L0) at (0,0) {$L_0$};
  \node[vtx] (A0) at (2.1,0) {$A_0$};
  \node[vtx] (D1) at (4.2,0) {$D_1$};
  \node[vtx] (L1) at (6.3,0) {$L_1$};
  \node[vtx] (A1) at (8.4,0) {$A_1$};
  \node[vtx] (D2) at (10.5,0) {$D_2$};
  \node[lat] (U0) at (1.0,-1.9) {$U_0$};
  \node[lat] (U1) at (7.3,-1.9) {$U_1$};
  \draw[ed] (L0) -- (A0);
  \draw[ed] (A0) -- (D1);
  \draw[ed] (D1) -- (L1);
  \draw[ed] (L1) -- (A1);
  \draw[ed] (A1) -- (D2);
  \draw[ed] (L0) to[bend left=18] (D1);
  \draw[ed] (L0) to[bend left=30] (L1);
  \draw[ed] (A0) to[bend left=18] (L1);
  \draw[ed] (A0) to[bend left=28] (A1);
  \draw[ed] (D1) to[bend left=18] (A1);
  \draw[ed] (L1) to[bend left=18] (D2);
  \draw[edd] (U0) -- (L0);
  \draw[edd] (U0) -- (A0);
  \draw[edd] (U0) to[bend right=12] (D1);
  \draw[edd] (U1) -- (L1);
  \draw[edd] (U1) -- (A1);
  \draw[edd] (U1) to[bend left=14] (D2);
  \draw[bid] (U0) -- (U1);
  \foreach \n/\x/\y/\t in {L0/0/0/{$L_0$},A0/2.1/0/{$A_0$},D1/4.2/0/{$D_1$},L1/6.3/0/{$L_1$},A1/8.4/0/{$A_1$},D2/10.5/0/{$D_2$}} {
    \node[vtx] at (\x,\y) {\t};}
  \node[lat] at (1.0,-1.9) {$U_0$};
  \node[lat] at (7.3,-1.9) {$U_1$};
\end{tikzpicture}
\caption{}
\label{fig:21.4}
\end{figure}

\section{大 g 公式}

本章与前两章偏重依赖于给定已测协变量 $L$ 的序贯可交换性、并以 g 公式做识别的方法。原因在于，实践中很少有复杂纵向数据的因果分析依赖其他识别条件与公式。例如，基于“前门公式是识别公式”那类识别条件的现实应用很少。然而，无论实质可信性与实际应用如何，不同的识别条件及其公式，在数学上都与序贯可交换性、以及基于所有变量（已测与未测）的 g 公式相联系，下面我们加以说明。

当给定已测协变量 $\bar L$ 序贯可交换性成立时，我们已经讨论了基于已测时变协变量 $\bar L$ 的 g 公式如何识别时变处理 $\bar A$ 对结局 $Y$ 的因果效应。现在假设我们有一个同时含观测变量 $(\bar A,\bar L,Y)$ 与未观测变量 $U$ 的因果 DAG，且已测变量 $\bar L$ 不足以实现序贯可交换性。

对任何因果 DAG，已测与未测变量的组合 $\bar X=(\bar L,U)$ 保证（联合）序贯可交换性，因为处理变量的任何父节点都包含在 $\bar A$ 或 $\bar X$ 中。因此，若因果图上的每个变量都被测量且正值性成立，基于 $\bar X$ 的 g 公式就能识别任何处理策略 $g$ 下的反事实均值 $\E[Y^g]$。我们把用 $\bar X$ 替换 $\bar L$ 的 g 公式称为\emph{大 g 公式}，因为它并不只基于观测数据。

\mnote{我们把 $(\bar A,\bar L,Y,U)$ 称为事实量（factuals），以区别于反事实量。事实量是实际世界中存在的变量。与观测变量不同，某些事实量（如 $U$）无法用于数据分析，通常是因为未被测量。}

给定因果 DAG、处理 $\bar A$ 与结局 $Y$、处理策略 $g$ 以及事实量 $(\bar A,\bar L,Y,U)$，我们可以显式写出 $Y^g$ 的分布（密度）的大 g 公式。大 g 公式只依赖于事实量 $(\bar A,\bar L,Y,U)$ 的分布。

大 g 公式是识别任何处理策略下反事实密度的正确公式，但大 g 公式不能在实践中使用，因为它含未测变量。于是有一个有趣的数学问题：大 g 公式能否化为观测数据 $(\bar A,\bar L,Y)$ 联合分布的泛函？若能，我们就得到一个不以 g 公式形式表达的新公式，它 (i) 重现大 g 公式的结果（因而是正确公式），且 (ii) 只以观测变量的分布写出（因而是可用于数据分析的公式）。

例如，在称为前门准则的识别条件下，$\E[Y^a]$ 的大 g 公式化为只含观测变量的公式——前门公式（证明见技术注记 21.11）。因此，在图 7.14 因果图所蕴含的前门假设下，前门公式是 $\E[Y^a]$ 均值的有效公式。

更一般地，我们希望回答以下两个问题。第一，我们是否总能判定：在除“$(\bar A,\bar L,Y,U)$ 的联合分布服从因果 DAG 所蕴含的 d 分离关系”之外不作任何假设时，大 g 公式能否改写为只依赖观测变量 $(\bar A,\bar L,Y)$ 分布的公式？第二，当前一问题的答案为是时，我们能否显式给出这样的识别公式？这两个问题都已得到肯定回答。

\mnote{这些问题已由 Tian 与 Pearl (2002)、Shpitser 与 Pearl (2006)、Huang 与 Valtorta (2006) 的工作完全解决。}

重要的是，这些是关于 $(\bar A,\bar L,Y,U)$ 上分布性质的纯数学问题，已知这些分布服从由 DAG 上 d 分离刻画的某些独立关系。也就是说，这些问题既不涉及反事实，也不涉及因果性。与因果性的唯一联系是“该 DAG 是因果 DAG”这一断言。若是，大 g 公式就有因果解释；若不是，肯定回答虽然仍为真，却没有因果含义。当然，在观察性分析中，我们永远无法确知我们猜想为因果图的图确实是因果图。

\begin{techpoint}[前门公式的大 g 公式证明]
在技术注记 7.4 中，我们给出了图 7.14 因果图下反事实概率 $\Pr[Y^a=y]$ 的前门公式的一个证明。这里我们用大 g 公式给出另一个证明。第二个证明依赖图 7.14 所蕴含的条件独立关系，但不要求反事实 $Y^m$ 存在。

图 7.14 下 $\Pr[Y^a=y]$ 的大 g 公式为
\[
\sum_m\sum_a\Pr[Y=y|M=m,A=a,U=u]\Pr[M=m|A=a,U=u]\Pr[U=u].
\]
由于没有 $U$ 的数据，$\Pr[Y^a=y]$ 可识别当且仅当大 g 公式只依赖于观测数据 $(Y,M,A)$ 的分布。下面我们说明确实如此，因为在上面的假设下，该 g 公式化为前门公式。

利用 d 分离，可以把大 g 公式改写为
\[
\begin{aligned}
&\sum_m\Pr[M=m|A=a]\sum_a\sum_u\Pr[Y=y|M=m,U=u]\\
&\qquad\times\Big\{\sum_{a'}\Pr[U=u|A=a']\Pr[A=a']\Big\}\\
&\quad=\sum_m\Pr[M=m|A=a]\sum_{a'}\sum_u\Pr[Y=y|M=m,A=a',U=u]\\
&\qquad\times\Pr[U=u|M=m,A=a']\Pr[A=a']\\
&\quad=\sum_m\Pr[M=m|A=a]\sum_{a'}\Pr[Y=y|M=m,A=a']\Pr[A=a'],
\end{aligned}
\]
最后一行用到 $U\indep M|A$ 与 $A\indep Y|M,U$，这正是前门公式。

我们再给出前门公式的另一个证明，它同样不要求反事实 $Y^m$ 存在。在确立 $\Pr[Y^a=y]$ 是由大 g 公式给出的 $(Y,M,A,U)$ 分布的函数之后，我们可以使用耦合论证。假设所有人都在实质层面同意，定义良好的 $Y^m$ 并不存在。然而，任何关于图 7.14 满足马尔可夫性的事实分布，都与一个“与数据同样细致”的底层 FFRCISTG 模型相容（Robins and Richardson, 2010），而按定义，该模型形式上包含变量 $Y^m$。技术注记 7.4 的证明表明，在该模型下大 g 公式等于前门公式。由此可知，不存在关于图 7.14 满足马尔可夫性、却使该等式失败的事实分布；因为若失败，该事实分布就不能与一个“与数据同样细致”的 FFRCISTG 模型相容。技术注记 21.12 给出基于 SWIG 性质的前门公式另一证明。
\end{techpoint}

\begin{techpoint}[用 SWIG 上处理节点的 d 分离证明前门公式]
这里我们用 SWIG 一个尚未讨论的重要性质，给出前门公式的另一个证明。

给定因果图 $G$，令 $G^a$ 为策略 $a$ 对应的 SWIG，$B^a$ 与 $C^a$ 为观测到的非处理节点 $(Y^a,\bar L^a)$ 的两个不相交子集。我们只假设处理反事实定义良好。SWIG $G^a$ 满足如下性质（Shpitser et al., 2022）：若固定节点 $a_m$ 在给定 $C^a$ 时与 $B^a$ 满足 d 分离，则 $\Pr(B^a=b|C^a=c)$ 不依赖于 $a_m$。这一性质与先前讨论的“任何以处理 $a_m$ 为内点的路径都被阻断”并不冲突。为说明新性质的含义，考虑图 7.14 前门图所蕴含的 SWIG $G^a$。在 SWIG $G^a$ 上，定义 $B^a=Y^a$、$C^a=(M^a,A)$。于是给定 $C^a$ 时 $a$ 与 $B^a$ 满足 d 分离，因为从 $a$ 到 $Y^a$ 的唯一路径经过 $C^a$ 中的非对撞节点 $M^{a'}$。因此按上述性质，对任意 $a$ 与 $a'$ 都有 $\E[Y^a|M^a,A]=\E\big[Y^{a'}|M^{a'},A\big]$。注意该性质不是跨世界的；它刻画的是不同单世界反事实分布之间的关系。

现在，在定义良好的反事实 $Y^m$ 不存在时，我们用这一 SWIG 性质证明前门公式。我们继续假设 $(Y^a,M^a,A)$ 按 SWIG $G^a$ 分解，而 $(Y^{a'},M^{a'},A)$ 按 SWIG $G^{a'}$ 分解。我们沿用技术注记 7.4 的证明，直到必须证明
\[
\E[Y^a|M^a]=\sum_{a'}\E[Y|M,A=a']\Pr(A=a').
\]
此时有 $\E[Y^a|M^a]=\sum_{a'}\E[Y^a|M^a,A=a']\Pr(A=a'|M^a)=\sum_{a'}\E[Y^a|M^a,A=a']\Pr(A=a')$，因为 $M(a)$ 与 $A$ 满足 d 分离。我们的新 SWIG 性质蕴含 $\E[Y^a|M^a,A=a']=\E\big[Y^{a'}|M^{a'},A=a'\big]=\E[Y|M,A=a']$，其中最后一个等式由一致性得到。于是 $\E[Y^a|M^a]=\sum_{a'}\E[Y|M,A=a']\Pr(A=a')$，如所要求。有趣的是，由此可知，尽管对所有 $a,a'$ 都有 $\E[Y^a|M^a]=\E\big[Y^{a'}|M^{a'}\big]$，却有 $\E[Y^a|M^a]\neq\E[Y|M]$，因为 $\E[Y|M]=\sum_{a'}\E[Y|M,A=a']\Pr(A=a'|M)$，而与反事实 $M^a$ 不同，观测到的事实 $M=M^A$ 并不独立于 $A$。
\end{techpoint}

\begin{techpoint}[一般结构嵌套均值模型的形式定义]
Robins (2004) 指出，在结构嵌套均值模型（SNMM）中，把“末次脉冲之后遵循的策略”取为 $\bar 0$ 并没有什么特别之处。我们可以改为相对于任意策略 $g$ 来定义脉冲函数，如下。给定 $g=(g_0,g_1,\dots,g_K)$，加性 SNMM 是这样一个模型：它刻画在时刻 $t$ 施加大小为 $a_t$ 的处理脉冲、随后从时刻 $t+1$ 起遵循 $g$，相对于从时刻 $t$ 起就遵循 $g$，对 $Y$ 的因果效应（以截至时刻 $t$ 的处理与协变量历史为条件）。也就是说，加性 SNMM 建模反事实对比
\[
\gamma^g_t(a_t,l_t)=\E\big[Y^{a_{t-1},a_t,g_{t+1}}-Y^{a_{t-1},g_t,g_{t+1}}\big|\bar A_{t-1}=\bar a_{t-1},A_t=a_t,\bar L_t=\bar l_t\big]
\]
其中 $t=0,\dots,K$，$a=(a_0,a_1,\dots,a_K)$，$g_{t+1}=(g_{t+1},\dots,g_K)$。当我们想强调 $a_t$ 与 $g_t$ 的独特作用时，把 $\gamma^g_t(a_t,l_t)$ 写作 $\gamma^g_t(\bar a_{t-1},a_t,\bar l_t)$，把 $Y^{a_{t-1},g_t}$ 写作 $Y^{a_{t-1},g_t,g_{t+1}}$。注意当 $a_t=g_t(\bar a_{t-1},\bar l_t)$ 时 $\gamma^g_t(a_t,l_t)\equiv0$。若如正文那样假设序贯可交换性，则可以从 $\gamma^g_t(a_t,l_t)$ 定义的条件事件中删去 $A_t=a_t$。

SNMM 假设 $\gamma^g_t(a_t,l_t)=\gamma^g_t(a_t,l_t;\beta)$，其中 $\gamma^g_t(a_t,l_t;\beta^\dagger)$ 是已知函数，当有限维参数向量 $\beta^\dagger$ 等于 0 或 $a_t=g_t(\bar a_{t-1},\bar l_t)$ 时取 0。若定义
\[
H_k(\gamma^g)=Y-\sum_{t=k}^{K}\gamma^g_t(A_t,L_t),
\]
则由一致性本身（Robins 2004）可得，对 $k=0,\dots,K$ 有 $\E[H_k(\gamma^g)|\bar L_k,\bar A_k]=\E[Y^{\bar A_{k-1},g_k}|\bar L_k,\bar A_k]$，且 $\E[H_0(\gamma^g)]=\E[Y^g]$。因此，若能识别 $\gamma^g_t(a_t,l_t)$，就能识别 $\E[Y^{\bar A_{k-1},g_k}|\bar L_k,\bar A_k]$ 与 $\E[Y^g]$。在正值性与序贯可交换性下，最后这个方程蕴含 $\E[H_k(\gamma^g)|\bar L_k,\bar A_k]=\E[H_k(\gamma^g)|\bar L_k,\bar A_{k-1}]$，这意味着 $\gamma^g_t(a_t,l_t)$ 可非参数识别。Robins (2004) 还定义了最优方案结构嵌套模型（opt-SNMM），并说明在正值性与序贯可交换性下，如何用 opt-SNMM 估计最优处理策略 $g_{opt}=\arg\max_g[\E(Y^g)]$。

但序贯可交换性并不是唯一可能的识别假设。例如，Zahn 等 (2022) 说明，$\gamma^g_t$ 可在一个时变平行趋势假设下识别，该假设推广了时变处理与协变量的双重差分估计通常所用的识别假设。

在技术注记 21.9 中，我们把 SNMM 中的 $g$ 取为“从不治疗”策略，即 $g=0$，并描述了在序贯可交换性假设下识别任意策略 $g$ 的 $\E[Y^g]$ 的算法。当序贯可交换性不成立时，我们可以用其他假设（例如时变平行趋势），它们足以识别用于定义 SNMM 的那个 $g$ 对应的 $\E[Y^g]$，却不足以识别任何其他策略 $g'$ 的 $\E[Y^{g'}]$。要识别后者，需要额外的假设。例如，Shahn 等 (2022) 说明，若在时变平行趋势之外还假设：在给定过去处理与已测协变量历史的条件下，不存在由未测混杂因子 $U$ 造成的加性效应修饰，则对所有 $g'$ 都能识别 $\E[Y^{g'}]$。这意味着最优策略 $g_{opt}=\arg\max_{g'}\E[Y^{g'}]$ 可识别。Shahn 等 (2022) 说明了如何用结构嵌套均值模型估计 $g_{opt}$。
\end{techpoint}
